\begin{thebibliography}{99}\setlength{\itemsep}{2pt}\small
\bibitem[Alfonsi et~al.(2012)]{alfonsi2012} Alfonsi, A., A. Schied, and A. Slynko (2012). Order book resilience, price manipulation, and the positive portfolio problem. \emph{SIAM Journal on Financial Mathematics} 3. doi:10.1137/110822098.
\bibitem[Alfonsi et~al.(2016)]{alfonsi2016} Alfonsi, A., A. Schied, and F. Kl\"ock (2016). Multivariate transient price impact and matrix-valued positive definite functions. \emph{Mathematics of Operations Research} 41, 914--934.
\bibitem[Almgren and Chriss(2001)]{almgren2001} Almgren, R., and N. Chriss (2001). Optimal execution of portfolio transactions. \emph{Journal of Risk} 3, 5--39.
\bibitem[Amilon(2008)]{amilon2008} Amilon, H. (2008). Estimation of an adaptive stock market model with heterogeneous agents. \emph{Journal of Empirical Finance} 15, 342--362.
\bibitem[Andrews and Cheng(2012)]{andrews2012} Andrews, D.\,W.\,K., and X. Cheng (2012). Estimation and inference with weak, semi-strong, and strong identification. \emph{Econometrica} 80, 2153--2211.

\bibitem[Banerjee and Szydlowski(2026)]{banerjee2026} Banerjee, S., and M. Szydlowski (2026). Trading against algorithms: Price dynamics and risk-sharing in a market with Q-learners. Working paper.
\bibitem[Barberis et~al.(2015)]{barberis2015} Barberis, N., R. Greenwood, L. Jin, and A. Shleifer (2015). X-CAPM: An extrapolative capital asset pricing model. \emph{Journal of Financial Economics} 115, 1--24.
\bibitem[Barberis et~al.(2018)]{barberis2018} Barberis, N., R. Greenwood, L. Jin, and A. Shleifer (2018). Extrapolation and bubbles. \emph{Journal of Financial Economics} 129, 203--227.
\bibitem[Barth and Kahn(2021)]{barth2021} Barth, D., and R.\,J. Kahn (2021). Hedge funds and the Treasury cash-futures disconnect. Office of Financial Research Working Paper 21-01.

\bibitem[Batenkov et~al.(2021)]{batenkov2021} Batenkov, D., G. Goldman, and Y. Yomdin (2021). Super-resolution of near-colliding point sources. \emph{Information and Inference} 10, 515--572.
\bibitem[Bernstein(1914)]{bernstein1914} Bernstein, S. (1914). Sur la meilleure approximation de $|x|$ par des polyn\^omes de degr\'es donn\'es. \emph{Acta Mathematica} 37, 1--57.
\bibitem[Bessembinder et~al.(2006)]{bessembinder2006} Bessembinder, H., W. Maxwell, and K. Venkataraman (2006). Market transparency, liquidity externalities, and institutional trading costs in corporate bonds. \emph{Journal of Financial Economics} 82, 251--288.
\bibitem[Biais et~al.(2015)]{biais2015} Biais, B., T. Foucault, and S. Moinas (2015). Equilibrium fast trading. \emph{Journal of Financial Economics} 116, 292--313.
\bibitem[Boehmer et~al.(2005)]{boehmer2005} Boehmer, E., G. Saar, and L. Yu (2005). Lifting the veil: An analysis of pre-trade transparency at the NYSE. \emph{Journal of Finance} 60, 783--815.
\bibitem[Bouchaud et~al.(2004)]{bouchaud2004} Bouchaud, J.-P., Y. Gefen, M. Potters, and M. Wyart (2004). Fluctuations and response in financial markets: The subtle nature of ``random'' price changes. \emph{Quantitative Finance} 4, 176--190.
\bibitem[Bouchaud et~al.(2018)]{bouchaud2018} Bouchaud, J.-P., J. Bonart, J. Donier, and M. Gould (2018). \emph{Trades, Quotes and Prices: Financial Markets Under the Microscope}. Cambridge University Press.
\bibitem[Brock and Hommes(1998)]{brock1998} Brock, W.~A., and C.~H. Hommes (1998). Heterogeneous beliefs and routes to chaos in a simple asset pricing model. \emph{Journal of Economic Dynamics and Control} 22, 1235--1274.
\bibitem[Brogaard et~al.(2014)]{brogaard2014} Brogaard, J., T. Hendershott, and R. Riordan (2014). High-frequency trading and price discovery. \emph{Review of Financial Studies} 27, 2267--2306.
\bibitem[Brunnermeier and Pedersen(2009)]{brunnermeier2009} Brunnermeier, M. K., and L. H. Pedersen (2009). Market liquidity and funding liquidity. \emph{Review of Financial Studies} 22, 2201--2238.
\bibitem[Castelli and Cover(1996)]{castelli1996} Castelli, V., and T.~M. Cover (1996). The relative value of labeled and unlabeled samples in pattern recognition with an unknown mixing parameter. \emph{IEEE Transactions on Information Theory} 42, 2102--2117.
\bibitem[CAT LLC(2026)]{catbudget2026} Consolidated Audit Trail, LLC (2026). Consolidated Audit Trail, LLC 2026 Financial and Operating Budget, 31 March 2026 (total expenses \$147{,}044{,}869). \url{https://www.catnmsplan.com/cat-financial-and-operating-budget}.
\bibitem[Cespa and Vives(2015)]{cespa2015} Cespa, G., and X. Vives (2015). The beauty contest and short-term trading. \emph{Journal of Finance} 70, 2099--2154.
\bibitem[Cochrane and Orcutt(1949)]{cochrane1949} Cochrane, D., and G. H. Orcutt (1949). Application of least squares regression to relationships containing auto-correlated error terms. \emph{Journal of the American Statistical Association} 44, 32--61.
\bibitem[Colliard et~al.(2022)]{colliard2022} Colliard, J.-E., T. Foucault, and S. Lovo (2022). Algorithmic pricing and liquidity in securities markets. CEPR Discussion Paper 17606.
\bibitem[Crouzet et~al.(2020)]{crouzet2020} Crouzet, N., I. Dew-Becker, and C.\,G. Nathanson (2020). On the effects of restricting short-term investment. \emph{Review of Financial Studies} 33, 1--43.
\bibitem[Dou et~al.(2025)]{dou2025} Dou, W.~W., I. Goldstein, and Y. Ji (2025). AI-powered trading, algorithmic collusion, and price efficiency. NBER Working Paper 34054.
\bibitem[Duarte and Young(2009)]{duarte2009} Duarte, J., and L. Young (2009). Why is PIN priced? \emph{Journal of Financial Economics} 91, 119--138.
\bibitem[Easley et~al.(1996)]{easley1996} Easley, D., N.~M. Kiefer, M. O'Hara, and J.~B. Paperman (1996). Liquidity, information, and infrequently traded stocks. \emph{Journal of Finance} 51, 1405--1436.
\bibitem[ESMA(2016)]{esma2016} European Securities and Markets Authority (2016). Guidelines on transaction reporting, order record keeping and clock synchronisation under MiFID II. ESMA/2016/1452.
\bibitem[European Commission(2017)]{rts24} European Commission (2017). Commission Delegated Regulation (EU) 2017/580 of 24 June 2016 supplementing Regulation (EU) No 600/2014 with regard to regulatory technical standards for the maintenance of relevant data relating to orders in financial instruments (RTS 24), Article 2(1)(b)--(c). \emph{Official Journal of the European Union} L 87, 31.3.2017.
\bibitem[European Commission(2017b)]{rts6} European Commission (2017). Commission Delegated Regulation (EU) 2017/589 of 19 July 2016 supplementing Directive 2014/65/EU with regard to regulatory technical standards specifying the organisational requirements of investment firms engaged in algorithmic trading (RTS 6). \emph{Official Journal of the European Union} L 87, 31.3.2017.
\bibitem[Fan(1991)]{fan1991} Fan, J. (1991). On the optimal rates of convergence for nonparametric deconvolution problems. \emph{Annals of Statistics} 19, 1257--1272.
\bibitem[Forssell and Ljung(1999)]{forssell1999} Forssell, U., and L. Ljung (1999). Closed-loop identification revisited. \emph{Automatica} 35, 1215--1241.
\bibitem[Franke and Westerhoff(2012)]{franke2012} Franke, R., and F. Westerhoff (2012). Structural stochastic volatility in asset pricing dynamics: Estimation and model contest. \emph{Journal of Economic Dynamics and Control} 36, 1193--1211.
\bibitem[Gatheral(2010)]{gatheral2010} Gatheral, J. (2010). No-dynamic-arbitrage and market impact. \emph{Quantitative Finance} 10, 749--759.
\bibitem[Glosten and Milgrom(1985)]{glosten1985} Glosten, L. R., and P. R. Milgrom (1985). Bid, ask and transaction prices in a specialist market with heterogeneously informed traders. \emph{Journal of Financial Economics} 14, 71--100.
\bibitem[Hasbrouck(1991)]{hasbrouck1991} Hasbrouck, J. (1991). Measuring the information content of stock trades. \emph{Journal of Finance} 46, 179--207.
\bibitem[Heinrich and Kahn(2018)]{heinrich2018} Heinrich, P., and J. Kahn (2018). Strong identifiability and optimal minimax rates for finite mixture estimation. \emph{Annals of Statistics} 46, 2844--2870.
\bibitem[Hendershott et~al.(2011)]{hendershott2011} Hendershott, T., C.~M. Jones, and A.~J. Menkveld (2011). Does algorithmic trading improve liquidity? \emph{Journal of Finance} 66, 1--33.
\bibitem[Hua and Sarkar(1990)]{hua1990} Hua, Y., and T. K. Sarkar (1990). Matrix pencil method for estimating parameters of exponentially damped/undamped sinusoids in noise. \emph{IEEE Transactions on Acoustics, Speech, and Signal Processing} 38, 814--824.
\bibitem[Katz et~al.(2023)]{katz2023} Katz, R., N. Diab, and D. Batenkov (2023). On the accuracy of Prony's method for recovery of exponential sums with closely spaced exponents. arXiv preprint arXiv:2302.05883.
\bibitem[Kirilenko et~al.(2017)]{kirilenko2017} Kirilenko, A., A. S. Kyle, M. Samadi, and T. Tuzun (2017). The Flash Crash: High-frequency trading in an electronic market. \emph{Journal of Finance} 72, 967--998.
\bibitem[Kong and Valiant(2017)]{kong2017} Kong, W., and G. Valiant (2017). Spectrum estimation from samples. \emph{Annals of Statistics} 45, 2218--2247.
\bibitem[Kreiss(1962)]{kreiss1962} Kreiss, H.-O. (1962). \"Uber die Stabilit\"atsdefinition f\"ur Differenzengleichungen die partielle Differentialgleichungen approximieren. \emph{BIT Numerical Mathematics} 2, 153--181.
\bibitem[Kukacka and Barunik(2017)]{kukacka2017} Kukacka, J., and J. Barunik (2017). Estimation of financial agent-based models with simulated maximum likelihood. \emph{Journal of Economic Dynamics and Control} 85, 21--45.
\bibitem[Kyle(1985)]{kyle1985} Kyle, A. S. (1985). Continuous auctions and insider trading. \emph{Econometrica} 53, 1315--1335.
\bibitem[Le Cam(1986)]{lecam1986} Le Cam, L. (1986). \emph{Asymptotic Methods in Statistical Decision Theory}. Springer.
\bibitem[Leach and Madhavan(1993)]{leach1993} Leach, J.\,C., and A.\,N. Madhavan (1993). Price experimentation and security market structure. \emph{Review of Financial Studies} 6, 375--404.

\bibitem[Liu and Ammari(2022)]{liu2023} Liu, P., and H. Ammari (2022). Super-resolution of positive near-colliding point sources. arXiv preprint arXiv:2212.00536.
\bibitem[Ljung(1999)]{ljung1999} Ljung, L. (1999). \emph{System Identification: Theory for the User}, 2nd ed. Prentice Hall, Upper Saddle River, NJ.
\bibitem[Makarov and Schoar(2020)]{makarov2020} Makarov, I., and A. Schoar (2020). Trading and arbitrage in cryptocurrency markets. \emph{Journal of Financial Economics} 135, 293--319.
\bibitem[Manole and Ho(2022)]{manole2022} Manole, T., and N. Ho (2022). Refined convergence rates for maximum likelihood estimation under finite mixture models. \emph{Proceedings of the 39th International Conference on Machine Learning}, PMLR 162.
\bibitem[Manski(2003)]{manski2003} Manski, C.~F. (2003). \emph{Partial Identification of Probability Distributions}. Springer, New York.
\bibitem[Moallemi et~al.(2008)]{moallemi2008} Moallemi, C.~C., B. Park, and B. Van Roy (2008). Strategic execution in the presence of an uninformed arbitrageur. arXiv preprint arXiv:0801.3001.
\bibitem[Moitra(2015)]{moitra2015} Moitra, A. (2015). Super-resolution, extremal functions and the condition number of Vandermonde matrices. In \emph{Proceedings of the 47th Annual ACM Symposium on Theory of Computing (STOC)}, 821--830.
\bibitem[Moskowitz et~al.(2012)]{moskowitz2012} Moskowitz, T. J., Y. H. Ooi, and L. H. Pedersen (2012). Time series momentum. \emph{Journal of Financial Economics} 104, 228--250.
\bibitem[Niroula et~al.(2026)]{niroula2026b} Niroula, P., S. Alkhushayni, Y. Chen, M. Bhavsar, and P. Niroula (2026). From human markets to machine markets: Transient amplification, an exact depth condition, and a maximal pre-trade guard. Working paper, Minnesota State University, Mankato.
\bibitem[Obizhaeva and Wang(2013)]{obizhaeva2013} Obizhaeva, A. A., and J. Wang (2013). Optimal trading strategy and supply/demand dynamics. \emph{Journal of Financial Markets} 16, 1--32.
\bibitem[Pagano and R\"oell(1996)]{pagano1996} Pagano, M., and A. R\"oell (1996). Transparency and liquidity: A comparison of auction and dealer markets with informed trading. \emph{Journal of Finance} 51, 579--611.
\bibitem[Park and Van Roy(2015)]{parkvanroy2015} Park, B., and B. Van Roy (2015). Adaptive execution: Exploration and learning of price impact. \emph{Operations Research}, doi:10.1287/opre.2015.1415.
\bibitem[Rebafka and Roueff(2015)]{rebafka2015} Rebafka, T., and F. Roueff (2015). Nonparametric estimation of the mixing density using polynomials. arXiv preprint arXiv:1002.4516.
\bibitem[Rivlin(1981)]{rivlin1981} Rivlin, T. J. (1981). \emph{An Introduction to the Approximation of Functions}. Dover.
\bibitem[Roll(1984)]{roll1984} Roll, R. (1984). A simple implicit measure of the effective bid-ask spread in an efficient market. \emph{Journal of Finance} 39, 1127--1139.
\bibitem[Roy and Kailath(1989)]{roy1989} Roy, R., and T. Kailath (1989). ESPRIT---estimation of signal parameters via rotational invariance techniques. \emph{IEEE Transactions on Acoustics, Speech, and Signal Processing} 37, 984--995.
\bibitem[SEC(2010)]{sec2010} U.S. Securities and Exchange Commission (2010). Risk management controls for brokers or dealers with market access. Release No.\ 34-63241.
\bibitem[SEC(2012)]{sec2012} U.S. Securities and Exchange Commission (2012). Consolidated Audit Trail; Final Rule (Release No.\ 34-67457, July 18, 2012, 77 FR 45722 at 45790--45795, pp.\ 220--238; proposing release No.\ 34-62174, May 26, 2010, 75 FR 32556 at 32598, 32600, pp.\ 147, 153).
\bibitem[SEC(2016)]{sec2016} U.S. Securities and Exchange Commission (2016). Joint Industry Plan; Order Approving the National Market System Plan Governing the Consolidated Audit Trail. Release No.\ 34-79318, November 15, 2016, 81 FR 84696 at 84801, p.\ 363.
\bibitem[SEC(2024)]{sec2024cat} U.S. Securities and Exchange Commission (2024). Notice of Filing of Proposed Rule Change Relating to the 2025 CAT Budget. Release No.\ 34-101901, December 12, 2024, p.\ 4 (revised May 2025).
\bibitem[SEC(2025)]{sec2025} U.S. Securities and Exchange Commission (2025). Order granting exemptive relief from certain requirements of the National Market System Plan governing the Consolidated Audit Trail. Release No.\ 34-102386, February 10, 2025.
\bibitem[Shahram and Milanfar(2005)]{shahram2005} Shahram, M., and P. Milanfar (2005). On the resolvability of sinusoids with nearby frequencies in the presence of noise. \emph{IEEE Transactions on Signal Processing} 53, 2579--2588.
\bibitem[Shleifer and Vishny(1992)]{shleifer1992} Shleifer, A., and R. W. Vishny (1992). Liquidation values and debt capacity: A market equilibrium approach. \emph{Journal of Finance} 47, 1343--1366.
\bibitem[Sornette et~al.(2023)]{sornette2023} Sornette, D., S. Lera, J. Lin, and K. Wu (2023). Non-normal interactions create socio-economic bubbles. \emph{Communications Physics} 6, 261.
\bibitem[Staiger and Stock(1997)]{staiger1997} Staiger, D., and J. H. Stock (1997). Instrumental variables regression with weak instruments. \emph{Econometrica} 65, 557--586.
\bibitem[Trebbi et~al.(2024)]{trebbi2024} Trebbi, F., M.~B. Zhang, and M. Simkovic (2024). The cost of regulatory compliance in the United States. Working paper.
\bibitem[Trefethen and Embree(2005)]{trefethen2005} Trefethen, L. N., and M. Embree (2005). \emph{Spectra and Pseudospectra: The Behavior of Nonnormal Matrices and Operators}. Princeton University Press.
\bibitem[Tremacoldi-Rossi(2025)]{tremacoldi2025} Tremacoldi-Rossi, P. (2025). Manipulating algorithmic markets. Working paper.
\bibitem[Tsybakov(2009)]{tsybakov2009} Tsybakov, A. B. (2009). \emph{Introduction to Nonparametric Estimation}. Springer.
\bibitem[van Waarde et~al.(2020)]{vanwaarde2020} van Waarde, H.~J., J. Eising, H.~L. Trentelman, and M.~K. Camlibel (2020). Data informativity: A new perspective on data-driven analysis and control. \emph{IEEE Transactions on Automatic Control} 65, 4753--4768.
\bibitem[Varga and Carpenter(1985)]{varga1985} Varga, R. S., and A. J. Carpenter (1985). On the Bernstein conjecture in approximation theory. \emph{Constructive Approximation} 1, 333--348.
\bibitem[Wagenmaker and Jamieson(2020)]{wagenmaker2020} Wagenmaker, A., and K. Jamieson (2020). Active learning for identification of linear dynamical systems. In \emph{Proceedings of the 33rd Conference on Learning Theory}, PMLR 125, 3487--3582.
\bibitem[Wagenmaker et~al.(2021)]{wagenmaker2021} Wagenmaker, A.~J., M. Simchowitz, and K. Jamieson (2021). Task-optimal exploration in linear dynamical systems. In \emph{Proceedings of the 38th International Conference on Machine Learning}, PMLR 139, 10641--10652.
\bibitem[Wellman(2006)]{wellman2006} Wellman, M.~P. (2006). Methods for empirical game-theoretic analysis. In \emph{Proceedings of the 21st National Conference on Artificial Intelligence (AAAI-06)}, 1552--1555.
\bibitem[Whittle(1953)]{whittle1953} Whittle, P. (1953). Estimation and information in stationary time series. \emph{Arkiv f\"or Matematik} 2, 423--434.
\bibitem[Wu and Yang(2020)]{wu2020} Wu, Y., and P. Yang (2020). Optimal estimation of Gaussian mixtures via denoised method of moments. \emph{Annals of Statistics} 48, 1981--2007.
\end{thebibliography}

\appendix
\section{Proofs}\label{app:proofs}

\paragraph{Lemma~\ref{lem:micro}.} (i) For $v_{t+1}=v_t+\zeta_{t+1}$ and $r_t=v_t+\nu_t$, the steady-state prior variance $P$ solves $P=P\sigma_\nu^2/(P+\sigma_\nu^2)+\sigma_\zeta^2$, i.e.\ $P^2-\sigma_\zeta^2P-\sigma_\zeta^2\sigma_\nu^2=0$, whose positive root is $P=\bigl(\sigma_\zeta^2+\sqrt{\sigma_\zeta^4+4\sigma_\zeta^2\sigma_\nu^2}\bigr)/2$. The gain is $\eta=P/(P+\sigma_\nu^2)\in(0,1)$, which simplifies to \eqref{eq:kgain}, and the update $\hat v_{t+1}=\hat v_t+\eta(r_t-\hat v_t)$ has the stated form. (ii) $\nabla_\theta\tfrac12H(x_t-\theta)^2=-H(x_t-\theta)$, so $\theta_{t+1}=(1-\gamma)\theta_t+\alpha H(x_t-\theta_t)=(1-\eta)\theta_t+\eta qx_t$ with $\eta=\gamma+\alpha H$ and $q=\alpha H/\eta$.

\paragraph{Lemma~\ref{lem:cm}.} Omit zero-weight cohorts. For $a>0$ put $b_i=\kappa q_i$ and
\[
T={\rm diag}\bigl(\sqrt a,\sqrt{b_1/\eta_1},\ldots\bigr),\qquad u=\bigl(\sqrt a,\sqrt{b_1\eta_1},\ldots\bigr)^\top .
\]
The price-first transition becomes the symmetric matrix $J'=TJT^{-1}={\rm diag}(0,r_1,\ldots)+uu^\top$, the input transforms to $\kappa u$ and $J'e_x=\sqrt a\,u$, so $h_j=(\kappa/a)e_x^\top(J')^{j+1}e_x$. The spectral decomposition of $J'$ gives nonnegative weights summing to $h_0=\kappa$, and Schur stability puts the eigenvalues in $[0,1)$. The case $a=0$ follows by continuity. A nonnegative mixture of decaying geometric sequences is completely monotone. Since $J$ is nonnegative with row sums at most one, every price row of $J^t$ has $\ell_1$, hence $\ell_2$, norm at most one, attained at $t=0$.

\paragraph{Theorem~\ref{thm:merged}.} With zero initial states, $y_{i,t}=\sum_{s\le t}\eta_ir_i^{t-s}x_s$, so $\sum_iq_iy_{i,t}=\sum_{s\le t}k_{t-s}x_s$ and every response is determined by $(a,\kappa,k)$, hence by $M$. Conversely, equal responses for all probes imply equal $H$, hence equal $G$, and partial fractions of $G$ identify the distinct poles with nonzero residues $w_i$, i.e.\ $M$. After cancellation, $H=\kappa zD(z)/[(z-a)D(z)-\kappa N(z)]$ with $D=\prod(z-r_i)$ and $N(z)=\sum_iw_iz\prod_{j\ne i}(z-r_j)$, as in \eqref{eq:H}; at a retained $r_i$ the denominator equals $-\kappa N(r_i)=-\kappa w_ir_i\prod_{j\ne i}(r_i-r_j)\neq0$ (memory rates are positive because $\eta_i<1$), so no root of $D$ cancels. The remaining numerator factor $z$ cancels if and only if the denominator vanishes at $0$; since $N(0)=0$ the denominator there equals $-aD(0)$, which is nonzero for $a>0$ and zero for $a=0$. Hence the McMillan degree, which is the dimension of a minimal realization, is $1+|{\rm supp}(M)\cap(0,1)|$ for $a'\neq0$ and $|{\rm supp}(M)\cap(0,1)|$ for $a'=0$ (an atom at rate $0$, a cohort with $\eta=1$, contributes the constant $M(\{0\})$ to $G$ and no pole, so it shifts $a$ to $a'=a+\kappa M(\{0\})$ and the argument above applies with $a'$) (numerically, the rank of the Hankel matrix of $(h_j)_{j\ge1}$ is $K+1$ and $K$ respectively for $K=1,2,3$; \texttt{hardening/p1\_checks.py}). For (ii), direct algebra gives
\[
G_\delta(z)-G_0(z)=-\frac{qz(z-1)\delta^2}{(z-1+\bar\eta)\bigl((z-1+\bar\eta)^2-\delta^2\bigr)} .
\]
On $|z|=1$ each factor has modulus at least $\epsilon$. Uniform stability gives $M_H=\sup|H_\theta|<\infty$, and the identity $H_\delta-H_0=z^{-1}H_\delta H_0(G_\delta-G_0)$, whose factor $z^{-1}$ has modulus one on the circle, gives an $\ell_2$ gain at most $C_{\mathcal F}\delta^2$. The mean difference of a power-constrained probe has squared norm at most $C_{\mathcal F}^2\delta^4PT$; the equal-noise KL is at most $C_{\mathcal F}^2\delta^4PT/(2\sigma_\nu^2)$, and Pinsker's inequality with total variation at least $1-2\alpha$ gives \eqref{eq:Tpower}.

\paragraph{Theorem~\ref{thm:quantity}.} In either positive economy $J\ge0$ and $\|J\|_\infty\le1-\epsilon(1-S)<1$, so the response is nonnegative with $\max h=\kappa$ and $\sum h=\kappa/(1-S)$. From the display above, $|G_\delta-G_0|\le2q\delta^2/\epsilon^3$ on the circle; combined with the resolvent identity and positivity this gives $\|H_\delta-H_0\|_\infty\le K_*\delta^2$, and Parseval bounds $\|h^\delta-h^0\|_2$ by the same constant. Young's inequality gives $\|\mu_\delta-\mu_0\|_2\le K_*\delta^2\|v\|_1$, and Gaussian KL with Pinsker gives \eqref{eq:klV}. For (ii), the mean difference is a lower-triangular convolution matrix times $v$; each column is a shifted copy of a prefix of $d_T$, so $\|\mu_\delta-\mu_0\|_2\le V\|d_T\|_2$ with equality for a pulse at $t=0$. For two known equal-covariance Gaussian laws the minimax error is $\Phi(-\|\mu_\delta-\mu_0\|_2/(2\sigma))$, attained by the midpoint likelihood-ratio test. For (iii), the merged measures have equal mass and their first moments differ by $-q\delta^2$, which enters $h_2$ with factor $\kappa^2$ while $h_0,h_1$ coincide; a pulse $A$ separates the third-observation means by $A\kappa^2q\delta^2$. Part (iv) follows from the $\ell_2$ gain bound, and (v) from superposition and $\max_jh_j=\kappa$, which give $\sup_t|\lambda x_t|\le|\lambda|\kappa V$, attained by a single pulse.

\paragraph{Lemma~\ref{lem:toeplitz}.} After the uncontrolled update, the pre-probe displacement at step $j$ is $\sum_{i<j}h_{j-i}v_i$; average linear-book execution adds $\kappa v_j/2$. Summing signed execution costs gives $\tfrac12v^\top T_hv$. Each kernel $\lambda^{|i-j|}$ is a positive-definite AR(1) covariance for $0\le\lambda<1$, and a nonnegative mixture with positive total weight is positive definite. For a round trip the fundamental cash term cancels.

\paragraph{Theorem~\ref{thm:cap}.} Let $U_t=\sum_{j<t}v_j$. Positive responses give $0\le x_t\le\kappa U_t$ and $I_t=-U_t$. Summing cash and marking inventory, $B-E_t=U_tx_t-\sum_{j<t}v_jx^0_{j+1}-\tfrac\kappa2\sum_{j<t}v_j^2+cU_t$. Its quadratic part $v^\top A_tv$ has diagonal $h_{t-1-i}-\kappa/2$ and off-diagonal $\tfrac12(h_{t-1-i}+h_{t-1-j}-h_{j-i})$ for $i<j$; monotonicity of $h$ bounds every entry by $\kappa/2$. Hence $B-E_t\le\kappa U_t^2/2+cU_t\le\kappa V^2/2+cV$, while position value is at most $V(P_*+\kappa V)$; the stated inequality guarantees margin. A single pulse attains both adverse bounds at its first response, so no larger allowance is robust. Solving the quadratic gives \eqref{eq:VB}. The quantity requirement $V_B\gtrsim\delta^{-2}$ then forces $B\gtrsim\delta^{-4}$, and the pulse $A$ funded with $(\mu+\tfrac12)\kappa A^2+(c+\mu P_*)A$ attains it.

\paragraph{Proposition~\ref{prop:disclosure}.} The initial cash receipt is common to both alternatives and the dealer then holds $-A$. At time three $E_{3,r}=B+\kappa A^2/2-cA-A^2h_{r,2}$, so $E_{3,\delta}-E_{3,0}=A^2\kappa^2q\delta^2>0$.

\paragraph{Lemma~\ref{thm:capital}.} Marking residual inventory at $P_*$ gives $E_\infty=B-\tfrac12v^\top T_hv-c\sum|v_t|\ge0$. For each mode $\lambda\le r$ the Toeplitz kernel has spectral density $(1-\lambda^2)/(1-2\lambda\cos\omega+\lambda^2)\ge(1-r)/(1+r)$, so $v^\top T_hv\ge m\|v\|_2^2$ and $\|v\|_2^2\le2B/m$. The transfer bound gives mean-difference squared norm at most $K_*^2\delta^4\|v\|_2^2$, hence the KL bound; Pinsker gives \eqref{eq:Bmin}. For one pulse, cash becomes $B_A-A(P_*+\kappa A/2)-cA$ and inventory $A$; at later marks $E_t=\mu AP_*+A^2h_{t-1}$ and $\mu|I_t|P_t=\mu AP_*+\mu A^2h_{t-1}$, so margin holds for $\mu\le1$.

\paragraph{Proposition~\ref{prop:tri}.} Parts (i) and (ii) follow from \eqref{eq:klconv} with $\Delta h=\Theta(\delta)$ for a rate shift and $\Theta(\delta^2)$ for symmetric splitting, and from additivity of KL over independent pulse cycles. Part (iii) solves the session constraint for $\delta$; substituting $d_T(\delta)\approx\delta^2\tilde d_T$ gives the fourth-root approximation.

\paragraph{Theorems~\ref{thm:spectral} and~\ref{thm:adaptive}.} For a flat round trip, telescoping the cash account gives $L_i(v)=\tfrac12v^\top T_{h_i}v=\frac1{2\pi}\int g_i(\omega)|V(e^{i\omega})|^2d\omega$, with $g_i>0$ by Lemma~\ref{lem:cm}. Parseval gives $\|(H_1-H_0)v\|_2^2=\frac1{2\pi}\int|\Delta\hat H|^2|V|^2\le\sup_\omega(|\Delta\hat H|^2/g_i)L_i(v)$, and the Neyman--Pearson requirement $\|(H_1-H_0)v\|_2\ge2\sigma z_{1-\alpha}$ gives \eqref{eq:Cdet}. For adaptive policies, Doob's optional stopping theorem applied to the log-likelihood-ratio martingale and the data-processing inequality give $\E_0[\Lambda_\tau]\ge{\rm kl}(\alpha,1-\beta)$, while $\E_0[\Lambda_\tau]=\E_0\sum_t(\Delta\mu_t)^2/(2\sigma^2)\le\sup(|\Delta\hat H|^2/g_0)\E_0[L_0]/(2\sigma^2)$, which gives \eqref{eq:adaptive}.

\emph{Part (iv).} Let $G_\theta$ be the $T\times T$ lower-triangular Toeplitz section of $H_\theta/\kappa$, so the observation covariance is $\Sigma_\theta=\sigma_w^2G_\theta G_\theta^\top+\sigma_\nu^2I$. A finite section of a bounded analytic symbol has norm at most its supremum on the circle, so $\|G_\theta\|\le\sup_\omega|H_\theta|/\kappa\le\bar M$ because $|H_\theta(e^{i\omega})|\le\sum_jh_j=\kappa/(1-S)$. The inverse of a lower-triangular Toeplitz section is the section of the reciprocal series; $\kappa/H_\theta$ is analytic on and outside the unit circle because $H_\theta$ is minimum phase (Theorem~\ref{thm:minphase}), and Lemma~\ref{lem:cm} gives ${\rm Re}\,H_\theta(e^{i\omega})=\sum_\ell\omega_\ell(1-\lambda_\ell\cos\omega)/(1-2\lambda_\ell\cos\omega+\lambda_\ell^2)\ge\sum_\ell\omega_\ell/(1+\lambda_\ell)\ge\kappa/2$, so $\|G_\theta^{-1}\|\le2$ and $\Sigma_\theta\succeq\sigma_{\rm eff}^2I$ (the smallest singular value of $G_\theta$ is at least $0.522$ in 762 random positive economies and sections; \texttt{hardening/p1\_checks.py}). The Gaussian divergence splits into a mean and a covariance part. The mean part is $\tfrac12\Delta\mu^\top\Sigma_1^{-1}\Delta\mu\le\|(H_1-H_0)v\|_2^2/(2\sigma_{\rm eff}^2)$. For the covariance part, the eigenvalues $\lambda_k$ of $A=\Sigma_1^{-1/2}\Sigma_0\Sigma_1^{-1/2}$ are at least $c=\sigma_\nu^2/(\sigma_\nu^2+\sigma_w^2\bar M^2)$, and $\lambda-1-\log\lambda\le(\lambda-1)^2/\lambda$, so it is at most $\|A-I\|_F^2/(2c)\le\|\Sigma_0-\Sigma_1\|_F^2/(2c\,\sigma_{\rm eff}^4)$. Writing $\Sigma_0-\Sigma_1=\sigma_w^2(\Delta G\,G_0^\top+G_1\Delta G^\top)$ and $\|\Delta G\|_F^2\le T\|h^\delta-h^0\|_2^2/\kappa^2\le TK_*^2\delta^4/\kappa^2$ (proof of Theorem~\ref{thm:quantity}) gives a covariance part of at most $C_wT\delta^4$. The data-processing inequality turns errors $(\alpha,\beta)$ into $\KL\ge{\rm kl}(\alpha,1-\beta)$, and part (ii) bounds $\|(H_1-H_0)v\|_2^2$ by $\sup_\omega(|\Delta\hat H|^2/g_i)L_i(v)$, which gives \eqref{eq:Cnoise}.

\paragraph{Theorem~\ref{thm:rh}.} Write a feasible flat probe as $v=Eu$ and expand $u=\sum_jc_je_j$ in the $G_i$-orthonormal eigenbasis, so $L_i(v)=\sum_jc_j^2$ and the separation constraint reads $\sum_j\lambda_jc_j^2\ge c_{\rm req}$. Write the top-two component as $(c_{L-2},c_{L-1})=r(\cos\theta,\sin\theta)$ with $r\ge0$; replacing $\theta$ by $\theta+\pi$ flips $d$ and $s_\theta$ and leaves \eqref{eq:rhdefs} unchanged, so $\theta\in[0,\pi)$ suffices. Put $L_0(\theta)=c_{\rm req}/\lambda(\theta)$ and $D=L_i(v)-L_0(\theta)$. The constraint gives $r^2\lambda(\theta)\ge c_{\rm req}-\sum_{j\le L-3}\lambda_jc_j^2$, hence
\[
D=r^2+\sum_{j\le L-3}c_j^2-L_0(\theta)\;\ge\;\sum_{j\le L-3}\Bigl(1-\frac{\lambda_j}{\lambda(\theta)}\Bigr)c_j^2\;\ge\;\Bigl(1-\frac{\lambda_{L-3}}{\lambda(\theta)}\Bigr)\sum_{j\le L-3}c_j^2\;\ge0,
\]
and $r^2\ge L_0(\theta)-\frac{\lambda_{L-3}}{\lambda(\theta)}\sum_{j\le L-3}c_j^2\ge L_0(\theta)\bigl(1-\gamma(\theta)D\bigr)$. By H\"older's inequality with the sign vector $s_\theta$, $c_{\rm fee}\|Eu\|_1\ge c_{\rm fee}r\|Ed(\theta)\|_1+c_{\rm fee}\sum_{j\le L-3}c_j\,s_\theta^\top Ee_j$, and Cauchy--Schwarz with weights $1-\lambda_j/\lambda(\theta)$ bounds the second term below by $-\sqrt{D\,W_\perp(\theta)}$. Using $r\ge\sqrt{L_0(1-\gamma D)_+}\ge\sqrt{L_0}\,(1-\gamma D)_+$ and the fact that the fee is nonnegative, with $x=\sqrt D$,
\[
L_i(v)+c_{\rm fee}\|v\|_1\;\ge\;L_0(\theta)+x^2+\Bigl[F_{\rm dom}(\theta)(1-\gamma(\theta)x^2)_+-\sqrt{W_\perp(\theta)}\,x\Bigr]_+\;\ge\;L_0(\theta)+\Psi(\theta).
\]
Minimizing over the unknown $\theta$ gives \eqref{eq:rh}, and $\Psi\ge0$, $L_0(\theta)\ge c_{\rm req}/\lambda_{L-1}$ give the last inequality. On the piece $\gamma x^2\le1$ with $\gamma F_{\rm dom}<1$ the bracket is a convex quadratic in $x$ with minimum $F_{\rm dom}-W_\perp/(4(1-\gamma F_{\rm dom}))$, which gives the closed form.

\paragraph{Corollary~\ref{cor:garch}.} Since $\Delta h_t=O(\delta^2)$, the variance difference $\sigma^2_{1,t}-\sigma^2_{0,t}=\sigma^2\xi[(h_1+h_0)*v]_t[\Delta h*v]_t=O(\delta^2)$, and the variance term of the Gaussian KL, $\tfrac12\sum(x-1-\log x)$ with $x=\sigma^2_{0,t}/\sigma^2_{1,t}$, is quadratic in that difference. Convex impact only increases dissipation.

\paragraph{Proposition~\ref{prop:crb}.} Standard Cram\'er--Rao for a Gaussian mean model with Jacobian $J$. For $K=2$ the columns of $J$ for the two rates become collinear as $\Delta\eta\to0$; expanding to third order, the smallest singular value of $J$ is of order $(\Delta\eta)^3$, so $[F^{-1}]_{kk}=\Theta((\Delta\eta)^{-6})$.

\paragraph{Theorem~\ref{thm:twocohort}.} Expanding $M_m(\theta_1)-M_m(\theta_0)$ under \eqref{eq:pair} gives zero for $m\le2$ and $4\tau p_-p_+(\Delta\eta)^3$ for $m=3$. The map from moments to impulse responses is analytic with nonsingular derivative at the common limit, so $\|\Delta h\|_2=\Theta(\tau(\Delta\eta)^3)$. Le Cam's two-point lemma with $N$ independent Gaussian experiments gives \eqref{eq:lecam2}. For the permutation-invariant loss, the identity matching of $\theta_0$ and $\theta_1$ costs $2\tau+(p_+-p_-)\Delta\eta=2\tau(1+\Delta\eta/q)$ per atom and the swapped matching costs $\max(2p_+,2p_-)\Delta\eta=(1+2\tau/q)\Delta\eta$; when both are at least $2\tau$, any estimate within $\tau$ of one economy in $d_\infty$ is farther than $\tau$ from the other by the triangle inequality, and the two-point argument is unchanged. For (iii), the upper bound is Theorem~\ref{prop:esprit}; the Prony--Hankel determinant identity follows from $g_m=w_1\eta_1r_1^m+w_2\eta_2r_2^m$.

\paragraph{Theorem~\ref{thm:kcohort}.} Let $(s^A,\omega^A)$ and $(s^B,\omega^B)$ be the $K$-point Gauss and Gauss--Radau rules for the uniform weight on $[-1,1]$, normalized to unit mass. Both are positive and exact for polynomials of degree $2K-2$, and their moments of degree $2K-1$ differ by $g_K\neq0$. With $r=r_c+\Delta\eta\,s$ and weights $W\omega$, the raw moments $\sum w_ir_i^m$ are polynomials in $\Delta\eta$ that agree below degree $2K-1$, so the kernels differ by $O((\Delta\eta)^{2K-1})$ with nonzero leading coefficient; analyticity of the kernel-to-response map gives $\|h^A-h^B\|_2\asymp(\Delta\eta)^{2K-1}$, and Le Cam's lemma gives the bound. For the $W_1$ version, the two merged measures are the images of the fixed normalized rules under $s\mapsto r_c+\Delta\eta\,s$ scaled to mass $W$, so their distance is $W\Delta\eta$ times the distance $2c_K>0$ between the rules, and the same two-point argument applies with the metric loss.

\paragraph{Lemma~\ref{lem:altpair}.} Let $L=2n-2$ and let $p^*$ be the best uniform approximation of $|x|$ on $[-1,1]$ by polynomials of degree $L$, with error $E=E_L(|x|)$. By Chebyshev's equioscillation theorem \citep{rivlin1981} there are points $t_0<\cdots<t_{L+1}$ at which $|x|-p^*(x)$ equals $\pm E$ with alternating signs. The divided-difference weights $w_i=1/\prod_{j\ne i}(t_i-t_j)$ alternate in sign and satisfy $\sum_iw_ip(t_i)=0$ for every polynomial $p$ of degree at most $L$; multiplying by $-1$ if necessary, $w_i$ has the sign of $|t_i|-p^*(t_i)$ for every $i$. Exactly $n$ weights are positive and $n$ negative, and the positive and negative parts have equal mass because constants are annihilated. Normalizing each part to unit mass and mapping $x\mapsto\rho(x+1)/2$ gives $n$-atom probability measures $\pi_A$ (positive part) and $\pi_B$ (negative part) on $[0,\rho]$ with equal moments of orders $0,\ldots,L$. For the 1-Lipschitz function $f(r)=|r-\rho/2|=(\rho/2)|x|$,
\[
\int f\,d(\pi_A-\pi_B)=\frac\rho2\int\bigl(|x|-p^*(x)\bigr)\,d(\pi_A-\pi_B)=\frac\rho2(E+E)=\rho E ,
\]
because $p^*$ integrates to the same value under both measures, $|x|-p^*=E$ on the support of $\pi_A$ and $-E$ on that of $\pi_B$. By Kantorovich--Rubinstein duality $W_1(\pi_A,\pi_B)\ge\rho E\ge c_*\rho/L$.

\paragraph{Theorem~\ref{thm:general}.} Write $K(z)=\sum_mk_mz^m$, $D=1-az-\kappa zK$ and $H=\kappa/D$, which is the generating function of the recursion $h_t=ah_{t-1}+\kappa\sum_mk_mh_{t-1-m}$ with $h_0=\kappa$. Since $k_m\ge0$ and $\sum_mk_m=\int(1-r)^{-1}d\mu\le W/(1-\rho)$, the $\ell_1$ norm of $az+\kappa zK$ is at most $s<1$, so $\|H\|_1\le\kappa/(1-s)$ and $\|1/H\|_1=\|D\|_1/\kappa\le(1+s)/\kappa$.

(i) Let $\mu_A=W\pi_A$ and $\mu_B=W\pi_B$ with $\pi_A,\pi_B$ from Lemma~\ref{lem:altpair}; both lie in $\mathcal M(\rho,W)$. With $\nu=\mu_A-\mu_B$, $\sum_m(k^A_m-k^B_m)z^m=\int(1-rz)^{-1}d\nu(r)=z^{2n-1}\int r^{2n-1}(1-rz)^{-1}d\nu(r)$, because the first $2n-1$ terms of the geometric series integrate to zero. On $|z|=1$ the modulus is at most $2W\rho^{2n-1}/(1-\rho)$, so Parseval gives $\|k^A-k^B\|_2\le2W\rho^{2n-1}/(1-\rho)$. Next, $H_A-H_B=z(K_A-K_B)H_AH_B$, so by Young's inequality $\|h^A-h^B\|_2\le\|H_A\|_1\|H_B\|_1\|k^A-k^B\|_2\le C_h\|k^A-k^B\|_2$. The Kullback--Leibler divergence between the $N$-sample Gaussian experiments is $N\|h^A-h^B\|_2^2/(2\sigma^2)\le(N/\sigma^2)\,2C_h^2W^2\rho^{2(2n-1)}/(1-\rho)^2$, which is at most one as soon as $2(2n-1)\log(1/\rho)\ge\Lambda_N-2\log(1/\rho)$. Take $n$ the smallest integer with this property and $n\ge2$; then $2n-2<\Lambda_N/(2\log(1/\rho))$ when $\Lambda_N>4\log(1/\rho)$. Since $d$ is a norm on signed measures and $d(\mu_A,\mu_B)=W\,W_1(\pi_A,\pi_B)\ge Wc_*\rho/(2n-2)>2c_*\rho W\log(1/\rho)/\Lambda_N$, the Bretagnolle--Huber form of Le Cam's lemma \citep[Lemma 2.6 and Theorem 2.2]{tsybakov2009} gives $\max_{j\in\{A,B\}}P_j(d(\hat\mu,\mu_j)\ge c_*\rho W\log(1/\rho)/\Lambda_N)\ge\tfrac14e^{-1}$.

(ii) \emph{Approximation.} Let $f$ be 1-Lipschitz on $[0,\rho]$ with $f(0)=0$, so $|f|\le\rho$. The function $g(x)=f(\rho(x+1)/2)$ is $(\rho/2)$-Lipschitz on $[-1,1]$, and Jackson's theorem \citep{rivlin1981} gives a polynomial $p$ of degree $L$ with $\sup|g-p|\le(\pi/2)(\rho/2)/(L+1)\le\rho$, hence $\sup|p|\le2\rho$. Its Chebyshev coefficients satisfy $|b_j|\le2\sup|p|\le4\rho$. The polynomial $q(r)=p(2r/\rho-1)$ approximates $f$ with the same error, and its monomial coefficients satisfy $\sum_j|q_j|\le\sum_{j\le L}|b_j|\,\|T_j(2r/\rho-1)\|_{\rm coef}\le4\rho\sum_{j\le L}B(\rho)^j\le\rho B(\rho)^{L+1}$, where $\|\cdot\|_{\rm coef}$ is the sum of absolute monomial coefficients, $\|T_j\|_{\rm coef}\le(1+\sqrt2)^j$, substituting $y=\alpha r-1$ multiplies the coefficient sum of $y^i$ by at most $(1+\alpha)^i$, and $B(\rho)-1\ge4$. On the event $\max_{j\le L}|\hat k_j-k_j|\le\delta$ the true $\mu$ attains objective at most $\delta$, so the minimizer $\hat\mu$ satisfies $\max_{j\le L}|\int r^jd\hat\mu-k_j|\le2\delta$ by the triangle inequality. \emph{Moment comparison.} If $\max_{j\le L}|\int r^jd\hat\mu-k_j|\le2\delta$, then $\hat\mu$ has mass at most $W+2\delta$ and
\[
\int f\,d(\hat\mu-\mu)=\int(f-q)\,d(\hat\mu-\mu)+\sum_{j\le L}q_j\Bigl(\int r^jd\hat\mu-k_j\Bigr)\le\frac{\pi\rho(W+\delta)}{2(L+1)}+2\delta\rho B(\rho)^{L+1}.
\]
\emph{Kernel recovery.} From $\kappa zK=1-az-\kappa/H$, the kernel values $k_0,\ldots,k_L$ are the coefficients of $(1-az-\kappa/H)/(\kappa z)$, which depend on $h_0,\ldots,h_{L+1}$ only. Work in the algebra of power series modulo $z^{L+2}$ with the $\ell_1$ norm, which is submultiplicative. If the averaged response is $\hat H=H+E$ and $\|1/H\|_1\|E\|_1\le\tfrac12$, then $1/\hat H=(1/H)(1+E/H)^{-1}$ and $\|1/\hat H-1/H\|_1\le2\|1/H\|_1^2\|E\|_1$, so $\max_j|\hat k_j-k_j|\le\delta\equiv2((1+s)/\kappa)^2(L+2)\max_{t\le L+1}|E_t|$. Each $E_t$ is Gaussian with variance $\sigma^2/N$, so with probability at least $1-\beta$, $\max_t|E_t|\le\sigma\sqrt{2\log(2(L+2)/\beta)/N}$. The true $\mu$ is feasible for the linear program with value at most $\delta$, so the minimizer satisfies the moment condition above. \emph{Choice of $L$.} With $L+1=\lfloor\log(N/\sigma^2)/(4\log B(\rho))\rfloor$, $B(\rho)^{L+1}\le(N/\sigma^2)^{1/4}$, so $\delta B(\rho)^{L+1}=O((\sigma^2/N)^{1/4}L\sqrt{\log(L/\beta)})$ and $\delta/(L+1)$ are $o(1/\log(N/\sigma^2))$, while $\pi\rho W/(2(L+1))\le2\pi\rho W\log B(\rho)/(\log(N/\sigma^2)-4\log B(\rho))$. For $N/\sigma^2$ larger than a threshold depending on $(\rho,W,a,\kappa,\beta)$ the side condition $\|1/H\|_1\|E\|_1\le\frac12$ holds on the event and the total is below $3\pi\rho W\log B(\rho)/\log(N/\sigma^2)$, uniformly over $\mathcal M(\rho,W)$. The final statement combines (i) and (ii).

\paragraph{Theorem~\ref{prop:passive}.} For stationary Gaussian processes with spectral densities $f_j\propto|H_j(e^{i\omega})|^2$, the divergence rate is
\[
D=\frac1{4\pi}\int_{-\pi}^{\pi}\bigl(\varrho-1-\log\varrho\bigr)\,d\omega,\qquad \varrho=f_1/f_0=|H_1/H_0|^2 .
\] By the resolvent identity, $H_1/H_0-1=z^{-1}H_1\Delta G$ with $\Delta G=G_1-G_0$, and uniform stability bounds $|H_1|$ on the circle, so $u\equiv\varrho-1=2\,{\rm Re}(z^{-1}H_1\Delta G)+O(|\Delta G|^2)$. Write $G(z)=\sum_kw_k\phi(z,\eta_k)$, with $w_k$ the cohort weights (any factor of $\eta_k$ is absorbed in $\phi$), where $\phi$ is analytic in $\eta$ on a neighbourhood of $[\eta_{\min},\eta_{\max}]$ uniformly for $|z|=1$ (its pole $z=1-\eta$ stays inside the disk). Expanding $\phi$ around $\bar\eta$, $\Delta G=\sum_{j\ge0}\frac{\partial_\eta^j\phi(z,\bar\eta)}{j!}\,\Delta\tilde M_j$, where $\tilde M_j=\sum_kw_k(\eta_k-\bar\eta)^j$. For the pair, $\Delta\tilde M_0=\Delta\tilde M_1=0$ and $\Delta\tilde M_2=q\delta^2$; for the moment-matched economies, $\Delta\tilde M_j=0$ for $j\le2$ and $\Delta\tilde M_3=4\tau p_-p_+(\Delta\eta)^3$ by \eqref{eq:dM3}; for the $K$-cohort Gauss and Gauss--Radau economies with rates $\bar\eta+\Delta\eta\,s$ and weights $W\omega$, both rules integrate polynomials of degree $2K-2$ exactly, so $\Delta\tilde M_j=W(\Delta\eta)^j\sum_k(\omega^A_k(s^A_k)^j-\omega^B_k(s^B_k)^j)$ vanishes for $j\le2K-2$ and equals $\pm Wg_K(\Delta\eta)^{2K-1}\neq0$ at $j=2K-1$. Higher terms are smaller by a factor of the spacing: $|\Delta\tilde M_j|\le2W(\Delta\eta)^j$, and analyticity of $\phi$ in $\eta$ on a disk of radius $\varrho>0$ around $\bar\eta$, uniformly on the circle, gives $|\partial_\eta^j\phi|/j!\le C\varrho^{-j}$, so for $\Delta\eta<\varrho/2$ the tail beyond the first nonzero term is $O((\Delta\eta)^{j_0+1})$. Hence $\sup_\omega|u|=O(\delta^2)$, $O(\tau(\Delta\eta)^3)$, respectively $O((\Delta\eta)^{2K-1})$, and $u=c(\omega)\epsilon+o(\epsilon)$ with $\epsilon=\delta^2$, $\tau(\Delta\eta)^3$ or $(\Delta\eta)^{2K-1}$ and $c(\omega)$ proportional to $2\,{\rm Re}\bigl(z^{-1}H_0\,\partial^{j_0}_\eta\phi(\cdot,\bar\eta)\bigr)/j_0!$, $j_0=2$, $3$ or $2K-1$. The function $c$ does not vanish identically: $f=z^{-1}H_0\partial^j_\eta\phi$ is analytic outside the disk including at infinity, where it vanishes, and is not identically zero; an analytic function on the exterior whose real part vanishes on the circle is an imaginary constant, which here would have to be zero. For $|u|\le\frac12$, $u^2/4\le u-\log(1+u)\le u^2$, so $D\asymp\int c(\omega)^2d\omega\,\epsilon^2$, which gives (i)--(iii). The sample requirements need a finite-$T$ bound that is uniform in the separation, which we obtain directly. The first $T$ observations are $N(0,\Sigma_j)$ with $\Sigma_j$ the $T\times T$ Toeplitz matrices of the spectral densities $f_j=\sigma^2|H_j(e^{i\omega})|^2$. On the compact, uniformly stable family the transfer functions have no zeros or poles on the circle (Theorem~\ref{thm:minphase}), so $0<\underline f\le f_j\le\bar f<\infty$ with constants that do not depend on the separation. A Toeplitz matrix has eigenvalues within the range of its symbol, so $\lambda_{\min}(\Sigma_0)\ge\underline f$ and $\|\Sigma_1-\Sigma_0\|\le\sup_\omega|f_1-f_0|=\sup_\omega f_0|u|\le\bar f\sup_\omega|u|$. With $A=\Sigma_0^{-1/2}\Sigma_1\Sigma_0^{-1/2}$, whose eigenvalues lie in $[\underline f/\bar f,\bar f/\underline f]$, the inequality $\lambda-1-\log\lambda\le(\lambda-1)^2/\lambda$ gives
\[
\KL(P_1^T\|P_0^T)=\tfrac12\sum_k(\lambda_k-1-\log\lambda_k)\le\frac{\bar f}{2\underline f}\,\|A-I\|_F^2\le\frac{\bar f}{2\underline f}\,T\Bigl(\frac{\bar f}{\underline f}\Bigr)^2\sup_\omega|u|^2\le C\,T\epsilon^2,
\]
with $C$ independent of $T$ and of the separation, because $\sup_\omega|u|=O(\epsilon)$ uniformly. The Bretagnolle--Huber inequality \citep[Lemma 2.6]{tsybakov2009} then gives the stated lower bounds on $T$ for every $T$, not only asymptotically. For a fixed pair, $T^{-1}\KL(P_1^T\|P_0^T)\to D$ and Stein's lemma (one error fixed) and the Chernoff bound (both errors vanishing) give the error exponents of the likelihood-ratio test, of the same order in the separation; that statement is pointwise in the pair.

\paragraph{Theorem~\ref{thm:minphase}.} In the delay variable $\lambda=z^{-1}$, $H=\kappa\prod_i(1-r_i\lambda)/\bigl[(1-a\lambda)\prod_i(1-r_i\lambda)-\kappa\lambda\sum_iw_i\prod_{j\ne i}(1-r_j\lambda)\bigr]$, which is \eqref{eq:H} divided through by $z^{K+1}$; its nonstructural zeros are the memory rates, which do not cancel under the nondegenerate assumptions of Theorem~\ref{thm:minphase}. In the z-plane there is also the structural timing zero at zero: the numerator is proportional to $z\prod_i(z-r_i)$. Stability puts the poles inside the disk. A rational spectral density $\sigma^2|H(e^{i\omega})|^2$ has a unique minimum-phase factor up to a unimodular constant, which is $H$; Theorem~\ref{thm:merged} then gives the merged measure.

\paragraph{Proposition~\ref{prop:tagged}.} With tagged flow, $u_{i,t}=r_iu_{i,t-1}+w_ix_t$ holds exactly and involves only cohort $i$, so $\tilde u_{i,t}=r_i\tilde u_{i,t-1}+w_ix_t+e_{i,t}$ with the MA(1) error $e_{i,t}$. Prices are independent of the measurement noise, so $\E[Z_te_{i,t}]=0$ and, by the ergodic theorem for the stationary linear price process, the IV estimator is consistent when $\Sigma_i=\E[Z_t(\tilde u_{i,t-1},x_t)]$ is nonsingular. Since $u_{i,t-1}=w_i\sum_{s\ge0}r_i^sx_{t-1-s}$, $\E[x_{t-1}u_{i,t-1}]=w_iA_0$ and $\E[x_tu_{i,t-1}]=w_iA_1$, which gives $\Sigma_i$. The price is a nonnegative mixture of geometric responses to its innovation (Lemma~\ref{lem:cm}), so $\gamma_x(s)=\sum_\ell c_\ell\lambda_\ell^{s}$ with $c_\ell=\sigma_\xi^2\kappa^{-2}\omega_\ell\sum_m\omega_m/(1-\lambda_\ell\lambda_m)\ge0$; substituting, $\det\Sigma_i=w_i(\gamma_x(0)A_0-\gamma_x(1)A_1)=w_i\sum_{\ell,m}c_\ell c_m(1-\lambda_\ell\lambda_m)/(1-r_i\lambda_m)$, each term positive. By summation by parts, $\sum_{t=1}^TZ_te_{i,t}=\sum_{t=1}^T(Z_t-r_iZ_{t+1})\varepsilon_{i,t}+O_P(1)$. Since $(\varepsilon_{i,t})$ is i.i.d.\ with mean zero and variance $\sigma_f^2$ and independent of the whole stationary price process, the summands $m_t=(Z_t-r_iZ_{t+1})\varepsilon_{i,t}$ form a stationary ergodic martingale difference sequence with respect to $\mathcal F_t=\sigma((x_s)_{s\in\mathbb Z},\varepsilon_{i,s}:s\le t)$, with $\E[m_tm_t^\top]=\sigma_f^2[(1+r_i^2)\Gamma_0-r_i(\Gamma_1+\Gamma_1^\top)]=\Omega_i$, so the Billingsley--Ibragimov martingale central limit theorem gives \eqref{eq:ivvar}. On a uniformly stable family the modes and weights lie in a compact set with $\lambda_\ell\le\bar\lambda<1$, so with $w_i\ge w_{\min}>0$ we get $\det\Sigma_i\ge w_{\min}(1-\bar\lambda^2)\gamma_x(0)^2>0$ and $\|\Omega_i\|$ is bounded, uniformly in the spacing.

\paragraph{Theorem~\ref{thm:blind}.} With $J_\theta$ upper triangular the eigenvalues are fixed. The stationary covariance solves $P=J_\theta PJ_\theta^\top+{\rm diag}(s_1,s_2)$, and $P_{11}$ is affine in $s_2$ with positive slope, so for each $\theta$ a unique $s_2(\theta)\ge0$ gives $P_{11}=V$ when $V$ exceeds the value at $s_2=0$. The one-step price variance is $s_1$ throughout, and $e_1^\top J_\theta e_2=\theta$ gives $\sup_k\|e_1^\top J_\theta^k\|\ge|\theta|$.

\paragraph{Proposition~\ref{prop:choice}.} (i) $\mathcal L_{\rm disc}<\mathcal L_{\rm price}$ is $F+Mc<M\ell G$. (ii) $T_p$ and $T_t$ are decreasing in $\Delta\eta$, with $T_t\ge\bar T$ exactly when $\Delta\eta\le\Delta\eta_t$ and $T_p\ge\bar T$ exactly when $\Delta\eta\le\Delta\eta_p$. On $\Delta\eta\le\Delta\eta_t$ both minima equal $\bar T$ and $G=0$. On $(\Delta\eta_t,\Delta\eta_p]$, $G=\bar T-c_t(\Delta\eta)^{-2}$, which increases to $G^*=\bar T-T_t(\Delta\eta_p)>0$. On $\Delta\eta>\Delta\eta_p$, $G=c_p(\Delta\eta)^{-(4K-2)}-c_t(\Delta\eta)^{-2}=(\Delta\eta)^{-2}\bigl[c_p(\Delta\eta)^{-(4K-4)}-c_t\bigr]$; the bracket is positive exactly when $\Delta\eta<\Delta\eta_0=(c_p/c_t)^{1/(4K-4)}$, and $G'(\Delta\eta)=(\Delta\eta)^{-3}\bigl[2c_t-(4K-2)c_p(\Delta\eta)^{-(4K-4)}\bigr]<0$ whenever $c_p(\Delta\eta)^{-(4K-4)}>c_t$, because $4K-2>2$. So $G$ decreases wherever it is positive beyond $\Delta\eta_p$, and is non-positive beyond $\Delta\eta_0$. Every superlevel set of $G$ at a positive level is therefore an interval (the negative tail, where $G$ rises back toward zero from below, never reaches a positive level), and the level $(F/M+c)/\ell>0$ lies above $G=0$, so the set where (i) holds is an interval contained in $(\Delta\eta_t,\Delta\eta_0)$ that contains $\Delta\eta_p$ whenever it is non-empty. (iii) The interval is non-empty if and only if the level is below the maximum, $F/M+c<\ell G^*$, that is $M>M^*$; the level decreases in $M$, so the superlevel set grows. $\square$

\section{Notation}\label{app:notation}

\begin{longtable}{p{0.18\textwidth}p{0.72\textwidth}}
\toprule
Symbol & Meaning\\
\midrule
\endhead
$x_t$ & pricing error (price minus fundamental)\\
$a$ & price persistence (anchor), $0\le a<1$\\
$\kappa$ & price impact per unit of executed quantity\\
$q_i,\ \eta_i,\ r_i$ & cohort weight, adaptation rate, memory rate $r_i=1-\eta_i$\\
$w_i=q_i\eta_i$ & merged weight\\
$M=\sum_iw_i\delta_{r_i}$ & merged rate--weight measure\\
$k_m,\ G(z)$ & feedback kernel and its transform\\
$H(z),\ h_j$ & closed-loop transfer function and impulse response, $h_0=\kappa$\\
$J,\ S$ & state transition matrix and $S=a+\kappa\sum q_i$\\
$v_t,\ V$ & executed probe quantity and total quantity $\sum|v_t|$\\
$d_T$ & impulse-response difference between hypotheses, first $T$ taps\\
$\sigma$ & observation-noise standard deviation\\
$\delta,\ \Delta\eta$ & half-separation of a known pair; spacing of unknown cohorts\\
$K_*$ & transfer-difference constant $2\kappa^2q/[\epsilon^3(1-S)^2]$\\
$T_h$ & Toeplitz matrix $(h_{|i-j|})$ of execution cost\\
$g_i(\omega)$ & dissipation kernel ${\rm Re}\,\hat H_i(e^{i\omega})-\kappa/2$\\
$B,\ \mu,\ c,\ \phi$ & capital, maintenance-margin fraction, proportional cost, taker fee\\
$V_{\rm eff}$ & per-trade execution capacity $\min(V_{\max},D_{\max}/\kappa)$\\
$\hat Q,\ \hat S$ & reduced-form loop gain and persistence indices\\
\bottomrule
\end{longtable}

\section{Internet appendix}\label{app:ia}

\subsection{A finite-window price-surveillance bound}

\begin{theorem}[Finite-window abstention]\label{thm:abstain}
Observe $x_t=sx_{t-1}+w_t$, $w_t\sim N(0,\sigma^2)$, $x_0=0$, $t\le T$, with $s\in(0,2)$. Let a possibly randomized rule $\Psi$ output Stable, Unstable or Abstain, with $\sup_{s<1}P_s(\Psi={\rm Unstable})\le\alpha$ and $\sup_{s\ge1}P_s(\Psi={\rm Stable})\le\alpha$, $\alpha<1/2$. Choose $0<c_0<1$ with $c_0e^{c_0}/(2\sqrt2)<1-2\alpha$. On $B_T=(1-c_0/T,1+c_0/T)$,
\[
\inf_{s\in B_T}P_s(\Psi={\rm Abstain})\;\ge\;1-2\alpha-\frac{c_0e^{c_0}}{2\sqrt2}>0 ,
\]
and for every finite $T$ the minimax worst-case error of binary rules over $s\in(0,2)$ is at least $1/2$.
\end{theorem}

\begin{proof}
Pair $s_0\in(1-c_0/T,1)$ with $s_1=s_0+c_0/T$. Then $\KL(P_{s_0}\|P_{s_1})=\frac{(s_1-s_0)^2}{2\sigma^2}\sum_t\E x_{t-1}^2\le c_0^2e^{2c_0}/4$, so Pinsker bounds total variation by $c_0e^{c_0}/(2\sqrt2)$. Error control gives $P_{s_0}({\rm Stable})\le\alpha+{\rm TV}$ and $P_{s_0}({\rm Unstable})\le\alpha$, which yields the abstention bound. For binary rules, continuity in total variation and Le Cam's two-point bound give worst-case error tending to $1/2$.
\end{proof}

The interval $B_T$ is an abstention layer that any error-controlled monitor must respect in the scalar Gaussian case; richer disclosure can remove it.

\subsection{Passive distinction of a known pair in finite samples}

\begin{proposition}[Passive distinction with zero controlled quantity]\label{prop:passive0}
Keep the known pair, but add independent uniform innovations $w_t\in[-b,b]$ of variance $s_w^2=b^2/3$ at each price update, observe $\tilde x_t=x_t+\nu_t$ with $\nu_t\sim N(0,\sigma^2)$, use independent three-step reset blocks and no probes, with $P_*>3b$. Let $b_1=a+\kappa q\eta$, $b_{2,0}=b_1^2+\kappa q\eta(1-\eta)$ and $b_{2,\delta}=b_{2,0}-\kappa q\delta^2$. In each block the covariance of the first and third observations differs by $\Delta=\kappa qs_w^2\delta^2$. With
\[
R=\max_{r\in\{0,\delta\}}\Bigl\{(s_w^2+\sigma^2)\bigl[s_w^2(1+b_1^2+b_{2,r}^2)+\sigma^2\bigr]+s_w^4b_{2,r}^2\Bigr\},
\]
the midpoint test of the average product has worst-case error at most $\alpha$ whenever $n\ge4R/(\alpha\Delta^2)$ blocks are observed: zero controlled quantity and order $\delta^{-4}$ blocks.
\end{proposition}

The Chebyshev bound is conservative but the order is real. With $b=0.03$, observation s.d.\ $0.01$ and $\delta=0.10$, the bound asks for $3.9\times10^7$ blocks at $\alpha=0.1$; in simulation the midpoint test still errs 32--38\% of the time with $2\times10^5$ blocks. Independent resets and known innovation laws are extra conditions. The proposition does not show that anonymous public records identify participants; it shows that passive excitation can accumulate memory information while the probe account stays untouched.

\begin{proof} In each block $x_1=w_1$, $x_2=b_1w_1+w_2$ and $x_3=b_{2,r}w_1+b_1w_2+w_3$, so $|x_t|\le3b$. Observation noise changes variances but not the covariance $s_w^2b_{2,r}$. The product variance is the product of marginal variances plus the squared covariance plus a fourth-cumulant term $b_{2,r}^2(\E w^4-3s_w^4)\le0$ for uniform innovations, so it is at most $R$. Chebyshev bounds each midpoint error by $4R/(n\Delta^2)$.
\end{proof}

\subsection{Blind inversion of the closed loop}\label{app:era}
\begin{enumerate}[itemsep=1pt]
\item \emph{Impulse decoupling:} $g_t=h_{t+1}-ah_t$.
\item \emph{Minimal realization:} form the Hankel matrix $H_{ij}=g_{i+j-1}$ and run Kung's eigensystem realization with a truncated SVD.
\item \emph{Transmission zeros:} compute the zeros of the realization, which under known $(\kappa,a)$ isolate the roots of the denominator of $G$.
\item \emph{Weights:} with recovered rates, project $u_t=(x_t-ax_{t-1})/\kappa$ on the reconstructed filters $y_{k,t-1}$.
\end{enumerate}
Without noise each step is exact: the realization has the order of the true system, its nonstructural transmission zeros are the memory rates (Theorem~\ref{thm:merged}); the known timing zero is discarded, and the projection returns the weights, so the procedure removes the pole-shift bias of fitting poles directly to the price response. With noise, per-cohort recovery degrades as rates move together at the rate given by Theorems~\ref{thm:kcohort} and~\ref{prop:esprit}, which is why we do not report it as an estimator of individual cohorts.

\subsection{Probing protocols}
Table~\ref{tab:protocols} compares four probing protocols at observation horizon $T=4$ per cycle.

\begin{table}[!htbp]\centering\small
\begin{threeparttable}
\caption{Probing protocols at observation horizon $T=4$ per cycle}\label{tab:protocols}
\begin{tabular}{lcrrrr}
\toprule
Protocol & $\delta$ & Executed vol. & Max trade & Error (SE) & Capital\\
\midrule
A: single pulse (unconstrained) & 0.03 & 131.5 & 131.5 & 0.047 (0.003) & 6{,}985.91\\
 & 0.05 & 47.3 & 47.3 & 0.054 (0.004) & 921.73\\
 & 0.08 & 18.5 & 18.5 & 0.053 (0.004) & 146.73\\
 & 0.12 & 8.2 & 8.2 & 0.044 (0.003) & 31.45\\
B: repeated capped pulses & 0.03 & 2{,}768.8 & 6.2 & 0.050 (0.003) & 7{,}035.75\\
 & 0.05 & 362.5 & 6.2 & 0.054 (0.004) & 923.88\\
 & 0.08 & 56.2 & 6.2 & 0.047 (0.003) & 146.00\\
 & 0.12 & 12.5 & 6.2 & 0.038 (0.003) & 34.88\\
C: sequential test, capped & 0.03 & 1{,}410.2 & 6.2 & 0.045 (0.003) & 3{,}585.14\\
 & 0.05 & 195.6 & 6.2 & 0.038 (0.003) & 500.00\\
 & 0.08 & 35.9 & 6.2 & 0.032 (0.003) & 94.30\\
 & 0.12 & 10.5 & 6.2 & 0.011 (0.002) & 29.85\\
D: generalized eigenvector & 0.03 & 2{,}359.7 & 3.7 & 0.050 (0.003) & 2{,}632.77\\
 & 0.05 & 314.8 & 3.7 & 0.048 (0.003) & 353.67\\
 & 0.08 & 49.2 & 3.7 & 0.046 (0.003) & 57.53\\
 & 0.12 & 9.9 & 3.8 & 0.045 (0.003) & 13.56\\
\bottomrule
\end{tabular}
\begin{tablenotes}\footnotesize\item 2{,}000 replications per hypothesis and cell. Log--log capital slopes: $-3.90$, $-3.84$, $-3.47$, $-3.81$. Protocols A and B reproduce the earlier version exactly. C and D are re-implemented because the earlier code was not retained: C is a Wald test with capped pulses; D repeats the top generalized eigenvector of a flat $(T+1)$-step probe, scaled to a maximum trade of 3.75, and is cheaper than the earlier D. Kind: \lab{recomputed} (\texttt{recompute\_p1\_protocols.py}).\end{tablenotes}
\end{threeparttable}
\end{table}

\subsection{Coordinate-weighted output}\label{app:coordw}
For $J=\begin{psmallmatrix}a&\kappa q\\ \eta a&1-\eta+\eta\kappa q\end{psmallmatrix}$ with $a=0.5$, $\kappa=0.8$, $q=0.2$, $\eta=0.3$ ($\rho(J)=0.8224$, non-normal), the price output has $\sup_t\|e_x^\top J^t\|=1$. The weighted output $C_w=(1,\lambda/{\rm MMR})=(1,160)$ with $\lambda=0.8$ and ${\rm MMR}=0.005$ has $\sup_t\|C_wJ^t\|=\|C_w\|=160.003$, attained at $t=0$, then $122.32$, $98.32$, $80.79$ \lab{recomputed in closed form}. The factor 160 is a static coordinate weight; $C_wz$ is not a margin quantity because $z$ contains no inventory. Margin must be computed on a cash-and-inventory ledger.

\subsection{Real-noise injection: requirements at each spacing}\label{app:injection}

Table~\ref{tab:injectionfull} gives the requirements behind Table~\ref{tab:injection}. $N^*$ is read off the grid $2^{j/2}$, so neighbouring values differ by a factor $\sqrt2$, and a ratio of $\sqrt2$ or $1/\sqrt2$ between the two arms is one grid step. Entries above $4{,}096$ use the central-limit draw described in Section~\ref{sec:injection}.

\begin{table}[!htbp]\centering\small
\begin{threeparttable}
\caption{Real-noise injection: number of experiments $N^*$ at each spacing}\label{tab:injectionfull}
\begin{tabular}{lrrrrr}
\toprule
 & $\Delta\eta$ & Price, real & Price, Gaussian & Tagged, real & Tagged, Gaussian\\
\midrule
$K=2$ & 0.500 & 64 & 64 & 724 & 1{,}024\\
 & 0.400 & 362 & 362 & 362 & 512\\
 & 0.320 & 1{,}448 & 1{,}024 & 362 & 362\\
 & 0.250 & 5{,}793 & 5{,}793 & 256 & 256\\
 & 0.200 & 16{,}384 & 23{,}170 & 256 & 362\\
 & 0.160 & 65{,}536 & 92{,}682 & 256 & 256\\
 & 0.125 & 370{,}728 & 370{,}728 & 256 & 256\\
\midrule
$K=3$ & 0.60 & 23{,}170 & 16{,}384 & & \\
 & 0.50 & 65{,}536 & 65{,}536 & & \\
 & 0.40 & 741{,}455 & 741{,}455 & & \\
 & 0.32 & 8{,}388{,}608 & 5{,}931{,}642 & & \\
 & 0.25 & 94{,}906{,}266 & 94{,}906{,}266 & & \\
\bottomrule
\end{tabular}
\begin{tablenotes}\footnotesize\item Real noise: 30{,}219 one-minute log-price changes of Binance BTCUSDT on 21 days (the 6th of each month, January 2025 to September 2026), kurtosis 19.4. Tagged requirements were computed for $K=2$ only. Kind: semi-synthetic.\end{tablenotes}
\end{threeparttable}
\end{table}

\subsection{Numerical stress tests of results that are not machine-checked}\label{app:stress}

Four results are proved on paper but not in Lean (Section~\ref{sec:mixtures} and Proposition~\ref{prop:tagged}). We stress-tested each across randomly drawn parameter settings. These are numerical checks, not proofs and not pre-specified tests; no prediction failed (Table~\ref{tab:stress}).

\begin{table}[!htbp]\centering\footnotesize
\begin{threeparttable}
\caption{Numerical stress tests (numerical checks, not proofs)}\label{tab:stress}
\setlength{\tabcolsep}{3pt}
\begin{tabular}{>{\raggedright\arraybackslash}p{0.22\textwidth}>{\raggedright\arraybackslash}p{0.33\textwidth}>{\raggedright\arraybackslash}p{0.37\textwidth}}
\toprule
Result & What was checked & Outcome\\
\midrule
Upper bound, Theorem~\ref{prop:esprit} & Prony's requirement $N^*$ for random separated configurations (15 for $K=2$, 12 for $K=3$); slope in $\Delta\eta$ and whether $N^*(\Delta\eta)^{4K-2}/\sigma^2$ grows as the spacing shrinks & Slopes $-6.72$ to $-5.79$ ($K=2$, theory $-6$) and $-11.01$ to $-9.29$ ($K=3$, theory $-10$); the normalized constant changes by a median factor of 1.05 and 0.94 from the widest to the narrowest spacing (at most 2.1); no configuration failed\\
Passive data, Theorem~\ref{prop:passive} & Exact Whittle divergence rate between two moment-matched economies, 10 random settings for each $K$ & Slopes $6.19$--$6.34$, $10.27$--$10.40$ and $14.33$--$14.54$ for $K=2,3,4$ (theory $6$, $10$, $14$): the divergence shrinks at least as fast as the theorem requires\\
General memory, Lemma~\ref{lem:altpair} & Alternation pairs on $[0,\rho]$ for 12 random $\rho\in[0.3,0.99]$ and $n=2,\ldots,6$: equal low moments, $W_1\ge\rho E_{2n-2}$, kernel distance bound & 60 of 60 cases pass\\
Value of disclosure, Theorem~\ref{thm:discvalue} & Tagged and price-only requirements over the spacing grid, $K=2,3,4$ & Tagged requirement does not grow as the spacing shrinks (it varies by factors 2.8, 11.3 and 5.7 over the grid); price-only slopes $-6.01$ and $-9.47$ for $K=2,3$; for $K=4$ our Prony implementation fails numerically at every spacing, so no price-only comparison\\
\bottomrule
\end{tabular}
\begin{tablenotes}\footnotesize\item Gaussian observation noise, $\sigma=10^{-3}$ (prices) and $\sigma_f=0.02$ (tagged flow), $a=0.5$, $\kappa=0.8$, total weight $0.08$, 200 replications per point, accuracy: every weight within 25\% of the smallest true weight in at least 75\% of replications. Moment-matched economies for the passive check: binomial weights on equally spaced nodes (the pair used in the Lean proof of Theorem~\ref{thm:kcohort}). Seed 20261031. Kind: numerical check.\end{tablenotes}
\end{threeparttable}
\end{table}

\section{Reproducibility: source of every number}\label{app:repro}

The replication package contains numbered Colab notebooks and stand-alone scripts; its README lists them in run order. ``Part $N$'' below and in table notes refers to notebook $N$ (Part 1: theorem checks and spot tests; Part 2: BTCUSDT futures depth; Parts 4--13: the cross-section, replications, fresh-day tests, second venues, the CFTC test, the certified guard and post-hoc checks). Result identifiers such as \texttt{X6} or \texttt{U6} name the JSON file that the notebook writes. ``Kind'' distinguishes closed-form or simulation results (recomputed), results from public market data (real data), and audits.

\begin{longtable}{p{0.39\textwidth}p{0.37\textwidth}p{0.14\textwidth}}
\toprule
Result & Source & Kind\\
\midrule
\endhead
Complete monotonicity, 892 economies & \texttt{L26\_complete\_monotonicity} & recomputed\\
$d_2=-\kappa^2q\delta^2$ & \texttt{P1.7-T1\_third\_tap} & recomputed\\
Ledger identity; $T_h$ positive definite & \texttt{R2-T1}, \texttt{R2-T3} & recomputed\\
Figure~\ref{fig:knownpair} & \texttt{make\_figs2.py} & recomputed\\
$\delta_{\min}(3)$, Table~\ref{tab:capbounds} recomputed columns, 61.14, 29.94 & \texttt{recompute\_theory.py} $\to$ \texttt{R\_P1\_theory\_constants.json} & recomputed\\
Pulse benchmarks (271, 58, 12, 11 pulses), Figures~\ref{fig:impulse}--\ref{fig:phase} & \texttt{make\_figs3.py} $\to$ \texttt{R\_P1\_pulse\_benchmarks.json} & recomputed\\
Table~\ref{tab:mcerr}, Table~\ref{tab:radau}, worked example, passive simulation & \texttt{recompute\_extra.py} $\to$ \texttt{R\_P1\_extra.json} & recomputed\\
Daily noise (Table~\ref{tab:roll}) & \texttt{P1.3-T1\_sigma\_obs} & real data\\
Table~\ref{tab:lecam} recomputed columns & same & recomputed\\
$K$-cohort slopes (Table~\ref{tab:kcohort}) & \texttt{T1\_K\_cohort\_rate} & recomputed\\
ESPRIT slope $-6.15$ & \texttt{P1.2-T1\_esprit\_rate} & recomputed\\
Minimum phase, 1,733 economies; rank 6 of 6 & \texttt{T2a}, \texttt{T2b} & recomputed\\
Passive slopes $-4.21$, $-5.93$ & \texttt{T3\_passive\_rate} & recomputed\\
Tagged-flow slope 0.011 & \texttt{T4\_disclosure} & recomputed\\
Spectral blindness gains & \texttt{T5\_spectral\_blindness} & recomputed\\
Noise and capital grid (Table~\ref{tab:capital}) & \texttt{P1.3-T1\_sigma\_obs} & real data\\
Loop gain and placebo (Table~\ref{tab:feedback}) & \texttt{P1.6-T1\_ema\_placebo} & real data\\
Estimator table (Table~\ref{tab:iv}) & \texttt{P1.5-T1}, \texttt{T2}, \texttt{T3} & recomputed\\
Data integrity (111 files) & \texttt{R2-T9\_data\_integrity} & audit\\
Tables~\ref{tab:ledger}, \ref{tab:capbounds} (certified, feasible, adaptive columns) & \texttt{recompute\_p1\_capital.py} $\to$ \texttt{R\_P1\_capital\_ledger.json} & recomputed\\
Table~\ref{tab:contracts}; resolution floors; Table~\ref{tab:protocols} & \texttt{recompute\_p1\_protocols.py} $\to$ \texttt{R\_P1\_protocols.json} & recomputed\\
SPRT capital figures & \texttt{recompute\_p1\_sprt.py} $\to$ \texttt{R\_P1\_sprt.json} & recomputed\\
Table~\ref{tab:crb} ($N^*$, $\sigma_{\rm crit}$) & \texttt{recompute\_crb.py} $\to$ \texttt{R\_P1\_crb.json} & recomputed\\
Table~\ref{tab:crb} (ML columns), Table~\ref{tab:lecam} (estimators) & \texttt{recompute\_p1\_estimators.py} & recomputed\\
Symmetric multi-asset network (Section~\ref{sec:blind}) & \texttt{recompute\_oct20.py} & recomputed\\
Proof audit: $K_*$ bound and $d_2$ identity (Theorem~\ref{thm:quantity}), moment gaps (Theorem~\ref{thm:twocohort}), spectral-blindness example (Theorem~\ref{thm:blind}) & \texttt{recompute\_audit.py} $\to$ \texttt{R\_audit\_oct2026.json} & recomputed\\
Theorem~\ref{prop:esprit}, Table~\ref{tab:kupper} & \texttt{kcohort\_upper.py} $\to$ \texttt{R\_P1\_kcohort\_upper.json} & recomputed\\
Lem.~\ref{lem:altpair}, Thm~\ref{thm:general}, Tab.~\ref{tab:general} & \texttt{general\_minimax.py} & recomputed\\
Theorem~\ref{thm:discvalue} numbers & \texttt{disclosure\_value.py} $\to$ \texttt{R\_P1\_disclosure\_value.json} & recomputed\\
Tables~\ref{tab:panel}, \ref{tab:assets}, \ref{tab:fee2} & Part 8: \texttt{V2\_spot\_panel}, \texttt{V3\_fee\_events} & real data\\
Table~\ref{tab:disc}, top panel & Part 4: \texttt{X9\_disclosure\_metrics} & real data\\
Table~\ref{tab:disc}, bottom panel; class-position forecasts & Part 6: \texttt{Z6\_C\_identification}, \texttt{Z6\_secondary}; \texttt{P6\_fresh\_days.json} (65acda36\ldots) & real data\\
Table~\ref{tab:cot} & Part 11: \texttt{C3\_cot\_data}, \texttt{C4\_cot\_disclosure}; \texttt{P11\_cot.json} (6f1d6f35\ldots) & real data\\
Basis-trade signature; real-data Cochrane--Orcutt & Part 13: \texttt{R6}, \texttt{R7} & real data (post hoc)\\
Feedback across 30 perpetuals & Part 4: \texttt{X3\_T6\_cross\_section} & real data\\
Theorem~\ref{thm:rh} full bound (Table~\ref{tab:rh}), direct minimization & \texttt{hardening/rh\_bound.py}, \texttt{rh\_check.py} $\to$ \texttt{R\_P1\_rh\_explicit.json} & recomputed\\
Minimal dimension at $a=0$; Toeplitz sections; passive slopes for $K\le4$; IV determinant & \texttt{hardening/p1\_checks.py} $\to$ \texttt{R\_P1\_hardening\_checks.json} & recomputed\\
Stylized nonlinear and trained online rules (Table~\ref{tab:learners}, Figure~\ref{fig:learners}) & \texttt{recompute\_\allowbreak p1\_\allowbreak learners.py} $\to$ \texttt{R\_P1\_\allowbreak learners.json} (Panel A); \texttt{validate\_\allowbreak trained\_\allowbreak learners\_\allowbreak hardened.py} $\to$ \texttt{R\_P1\_\allowbreak trained\_\allowbreak learners\_\allowbreak hardened.json} (Panels B--C); specification \texttt{P1\_\allowbreak trained\_\allowbreak learners.json} (3e8bf30c\ldots) & recomputed\\
Disclosure in dollars (Table~\ref{tab:cat_dollars}) & \texttt{recompute\_\allowbreak p1\_\allowbreak cat\_\allowbreak costs.py} $\to$ \texttt{R\_P1\_\allowbreak cat\_\allowbreak disclosure\_\allowbreak dollars.json} & recomputed\\
Out-of-sample post-shock return test (Section~\ref{sec:disc_data}) & \texttt{recompute\_\allowbreak prereg\_\allowbreak oct21\_\allowbreak 22nd.py} $\to$ \texttt{R\_P1\_\allowbreak disclosure\_\allowbreak oos.json}; \texttt{P1\_\allowbreak disclosure\_\allowbreak oos.json} (f0182ce5\ldots) & real data\\
Out-of-sample cohort-state trajectory test (Table~\ref{tab:disc_state_oos}) & \texttt{recompute\_\allowbreak prereg\_\allowbreak oct21\_\allowbreak followup.py} $\to$ \texttt{R\_P1\_\allowbreak disclosure\_\allowbreak state\_\allowbreak oos.json}; \texttt{P1\_\allowbreak disclosure\_\allowbreak state\_\allowbreak oos.json} (5136cc6b\ldots) & real data\\
\bottomrule
Proposition~\ref{prop:target}, Table~\ref{tab:target} & \texttt{theory/targeted\_intervention.py} $\to$ \texttt{R\_P1\_targeted\_intervention.json} & recomputed\\
Table~\ref{tab:lecam}, column $N^*_{\rm CRB}$; Table~\ref{tab:crb}, spacing-0.10 three-cohort row & \texttt{breakthrough\_experiments/results\_exp17\_minimax\_lecam\_bound.json}, \texttt{results\_exp15\_cohort\_information.json} & recomputed (earlier run)\\
Table~\ref{tab:rh}, best feasible probes and $L=320,480$ & \texttt{R\_P1\_capital\_ledger.json}, \texttt{R\_P1\_cquadL.json} (reproduced by \texttt{hardening/rh\_quad.py}) & recomputed\\
Table~\ref{tab:ledger}, peak margins of the eigenvector probes & \texttt{breakthrough\_experiments/results\_exp14\_capital\_frontier\_exact.json} & recomputed (earlier run)\\
Table~\ref{tab:general} & \texttt{theory/R\_P1\_general\_minimax.json} & recomputed\\
Table~\ref{tab:lecam} and Table~\ref{tab:crb}, estimator columns & \texttt{R\_P1\_estimators.json}, \texttt{R\_P1\_estimators\_crb.json} & recomputed\\
Theorem~\ref{thm:merged}, numerical check of the equivalence & \texttt{R2-T5\_merged\_invariance} & recomputed\\
Section~\ref{sec:disc_data} and Table~\ref{tab:scorecard}: horizon-matched test on 30 new markets (29 of 29, median 7.3, smallest 2.6) and jackknife-corrected persistence ($1.0008$) & \texttt{thesis/evidence/C2\_results.json} (specification \texttt{thesis/checks/prereg/C2\_margin\_vs\_identify.json}) & real data\\
Section~\ref{sec:disc_data}: exploratory 48- and 96-lag refit (median 7.5) & \texttt{thesis/evidence/pilot3\_horizon\_matched\_results.json} & real data (exploratory)\\
Section~\ref{sec:disc_data} and Table~\ref{tab:scorecard}: 105-day follow-up (median ratio 74.9; 9 identified; 10 of 30, 23 required) & \texttt{thesis/evidence/E1v2\_results.json} & real data\\
Section~\ref{sec:disc_data} and Table~\ref{tab:scorecard}: 58-day follow-up (22 of 29; 20 of 29, 22 required) & \texttt{thesis/evidence/D1\_results.json} & real data\\
Section~\ref{sec:injection}, Tables~\ref{tab:injection}, \ref{tab:injectionfull} and~\ref{tab:scorecard}: real-noise injection (slopes $-6.01$, $-9.81$; tagged 256--724) & \path{p1_injection/injection_test.py}, \path{p1_injection/results/I1_injection_test.json} (specification \path{p1_injection/prereg/P1_injection.json}, 596a700d\ldots) & semi-synthetic\\
Supplement~\ref{app:stress} and Table~\ref{tab:stress}: numerical stress tests & \path{p1_stress/stress_theorems.py}, \path{p1_stress/results/S1_stress_tests.json} & numerical check\\
Theorem~\ref{thm:merged}(i), Theorem~\ref{thm:kcohort} and Proposition~\ref{prop:tagged} (machine-checked parts) & \path{research_extension/current_manuscript_lean_20261008/CurrentProofs/Paper1MergedModel.lean}, \path{Paper1Identification.lean}, \path{Paper1MomentGap.lean}, \path{Paper1PoleCount.lean}, \path{Paper1FiniteIdentification.lean}, \path{Paper1Tagged.lean}, \path{Paper1Testing.lean}, \path{Paper1LowerBound.lean}, \path{Paper1MatchedPair.lean} & Lean 4 proof\\
Section~\ref{sec:cot}: post hoc robustness (futures prices, roll weeks, 26- and 52-week blocks) & \path{p11/robustness_cot.py}, \path{p11res/results/C4R_cot_robustness.json} & real data (post hoc)\\
Section~\ref{sec:cot} and Table~\ref{tab:scorecard}: class-flow test (3 of 8; $F\ge10$ in 7; $w\ne0$ in 6) & \path{p11/tagged_flow_test.py}, \path{p11res/results/C5F_cot_tagged_flow.json} (specification \path{p11/prereg_flow/P11F_cot_tagged_flow.json}, 052a63f2\ldots) & real data\\
\end{longtable}





% BEGIN ECONOMIC REVIEW 20261029
\section*{Development revision and validation status: 29 October 2026}
This is a working draft for author and specialist review. The revision corrects
hidden-state drift accounting, distinguishes cumulative displacement from peak
mark-to-market loss, and restricts claims to their stated model and observation
contracts. It does not declare every theorem independently verified or every
research objective complete.

\paragraph{Returned execution record.} The author returned a completed full-mode
notebook archive using source commit \texttt{19805d7e}, seed 20261029,
Python 3.13.18 and NumPy 2.1.3. All 10 exported-file hashes and the four analytical
source modules were checked. The run completed E1, E2, E3, E4 and E6; its log
reports 18 targeted tests passing. This is an execution reproduction of the same
seeded code, not independent scientific validation. No real-market CSV was
provided, and the GPU, known-source-feedback and prospective studies were not run.

\paragraph{What the checks establish.} E1 cross-checks the corrected drift identity
against a resolvent and exercises a sufficient invariant-box guard. E2 demonstrates
filter-induced persistence under martingale nulls. E3 compares known alternatives
with equal aggregate risk; it is an illustrative witness, not a uniform rate
theorem. E4 matches baseline capacity on training episodes and reports test
capacity differences; the guard and volatility control both have zero observed
budget breaches in the three bounded-shock synthetic regimes. This does not
establish universal guard superiority or economic welfare.

\paragraph{Coverage and the observation contract.} E6 uses 1,000 replications,
32 reset responses and 512 taps. The nominal 95\% interval assumes independent
Gaussian tap errors and a valid memory envelope. The intervals in the table
measure Monte Carlo uncertainty about coverage, not financial risk.
\begin{center}\small
\begin{tabular}{lrrc}
\toprule
Setting & Coverage & MC 95\% interval & Contract holds\\
\midrule
Independent Gaussian & 94.6\% & [93.0, 95.8] & yes\\
Correlated tap errors & 48.3\% & [45.2, 51.4] & no\\
Heavy-tailed readings & 95.1\% & [93.6, 96.3] & no\\
Biased readings & 25.6\% & [23.0, 28.4] & no\\
Unseen slow memory & 0.0\% & [0.0, 0.4] & no\\
\bottomrule
\end{tabular}
\end{center}
Failures under correlated errors, biased readings and a false memory floor expose
the contract's limits. The heavy-tail outcome in this configuration is not a
general robustness result. All outcomes, including negative ones, are retained.

\paragraph{Remaining obligations.} Independent review of the load-bearing
identification and certification proofs, continuous numerical extrema, economic
primitives, the public-observation bridge and policy counterfactuals remains
necessary. The current two-paper name-map audit finds 54 of 57 formal environments
with existing declaration mappings, three unmapped statements and seven revised
statements requiring semantic review. A label mapping is not a new Lean proof.
The experiment design and full results are retained under
\texttt{studies/economic\_review\_20261029/}.

% END ECONOMIC REVIEW 20261029

% BEGIN INFORMATION DECISION REVISION 20261029
\section*{Economic interpretation and information-to-decision extension}

\subsection*{What this revision adds}
The information question is distinct from choosing a preferred amount of risk.
For a specified financial decision, the relevant issue is which observations
support an exposure constraint, what remains unidentified, and whether obtaining
additional information changes the attainable decision after costs. This does
not introduce portfolio choice or control without full identification as new
ideas. The general identification-versus-decision-value result still requires
independent proof and a comparison with the closest existing results.

\subsection*{Time conventions and the institution}
The new application is a price-taking, unlevered own-account investor, not a
venue-wide supervisor. It begins with cash, buys once, holds a fixed position
for 60 one-minute model periods, and closes fully. Round-trip linear and
quadratic charges enter cash and wealth. The account cannot change the market's
feedback, and the experiment does not measure customer liquidity or welfare.
This is separate from the paper's dealer/venue admission contracts.

\begin{center}\small
\begin{tabular}{p{0.26\textwidth}p{0.65\textwidth}}
\toprule
Quantity & Interpretation\\
\midrule
Sampling interval & One minute in the new synthetic account; other historical studies retain their stated sampling clocks.\\
Historical lookback & Often six hours in the existing forecasting studies; not a holding period.\\
Forecast horizon & Often the next hour's adverse excursion; not an executed round-trip trade.\\
New holding period & Exactly 60 model minutes, with explicit entry and exit charges.\\
Other theory horizons & The existing finite-horizon theorems retain their declared $H$; no universal holding period is imposed.\\
\bottomrule
\end{tabular}
\end{center}

\subsection*{Information and decision design}
The new comparison uses controlled independent reset observations of 12 impulse
taps, a known plant and reading variance, and a declared finite candidate class.
It does not treat public market trading as a reset experiment. Common training
and validation paths supply 768 and 384 additional reset episodes per repetition;
512 further paths evaluate the selected actions. These are in addition to the
reported tap-reading budget, not evidence of learning from 32 observations in
total. Verified disclosure restricts
the minimum-memory rate to an audited band. The oracle knows the true response,
not future shocks. All certificate arms use the same training objective.

The account selects a position by empirical expected logarithmic terminal
wealth, subject to its information-supported exposure restriction. This is a
finite-horizon application of logarithmic utility, not a newly proved Kelly
growth-optimality theorem. Equal response bounds alone do not imply equal
expected returns or identical utility-optimal portfolios.

The new account has initial wealth of 1,000 dollars and an initial price of
100 dollars; its position grid runs from zero to 9.9 units. Its round-trip
dollar charge is $c(I)=0.002I+0.0002I^2$. For bounded external flow, the
constraint adds this charge to position times the largest admissible absolute
price response over the holding horizon. The resulting bound concerns loss
from initial wealth at each date. It is not peak-to-trough drawdown, protection
against arbitrary unbounded shocks, or a rolling deployment guarantee.

\subsection*{Returned execution and proof scope}
The returned full-mode archive uses seed 20261029. All 35 manifest entries and
eight core source files were checked. The run contains 20 independent
repetitions per regime, not 16,800 independent experiments: the 16,800 policy
rows reuse tapes across policies, budgets and observation counts. Twelve
targeted checks passed. Eight conditional decision-logic lemmas compiled with
Lean 4.19.0 and their axiom printouts report no axioms. This is a seeded
execution reproduction. It neither completes the older manuscript's Lean
coverage nor formalizes the Gaussian confidence argument, market dynamics,
general minimax separation, portfolio optimality or welfare.

The displayed slice uses 32 reset episodes and a one-dollar loss budget from
initial wealth of 1,000 dollars. It is an illustrative presentation of an
exploratory grid, not a primary selected before the results were seen. All
regimes, budgets, policies and negative outcomes remain in the result archive.
The inputs and complete outputs are retained with the verification record in
the new information-decision study directory.

\subsection*{Aggregate information and unresolved cohorts: returned simulation}
\begin{center}\small
\begin{tabular}{lrrr}
\toprule
Regime & Unique candidate (\%) & Certified units & Oracle units\\
\midrule
No feedback & 0.0 & 9.405 & 9.900\\
Moderate feedback & 5.0 & 5.000 & 5.000\\
Near boundary & 0.0 & 1.200 & 1.200\\
\bottomrule
\end{tabular}
\end{center}
Here ``unique'' means one remaining candidate in the declared finite class,
not identification of arbitrary real-world algorithms. In the moderate and
near-boundary configurations, the aggregate certificate reaches the oracle's
grid capacity in this slice even though the candidate decomposition is usually
unresolved. This is an illustrative decision comparison; it does not prove a
uniform rate separation for continuous classes.

\subsection*{Disclosure cost and the negative result}
\begin{center}\small
\begin{tabular}{lrrr}
\toprule
Paid minus aggregate & Mean (bp) & Lower (bp) & Upper (bp)\\
\midrule
No feedback & -0.100 & -0.100 & -0.100\\
Moderate feedback & -0.126 & -0.126 & -0.125\\
Near boundary & -0.100 & -0.100 & -0.100\\
\bottomrule
\end{tabular}
\end{center}
The intervals are exploratory percentile bootstraps over repetition-level
paired means, not simultaneous intervals across the grid. In this slice the
unpaid disclosure does not expand capacity, while the assumed one-cent charge
reduces net log wealth. This run therefore does not establish a positive net
economic value of disclosure. The price is hypothetical, not a calibrated
reporting cost. Stronger disclosure need not be worth purchasing when aggregate
information already supports the relevant decision.

\subsection*{Assumption failures and calibration diagnostics}
The table reports finite-sample frequencies for all seven regimes, combining
the three reading budgets. Each band outcome is counted once per repetition
and reading budget (60 outcomes per regime); repeating the same outcome across
four financial loss budgets would not create further observations. The three
reading budgets within a repetition are not an independent replication.
\begin{center}\small
\begin{tabular}{lcrr}
\toprule
Regime & Contract valid & Tap-band hit (\%) & Empty set (\%)\\
\midrule
No feedback & Yes & 96.7 & 3.3\\
Moderate feedback & Yes & 93.3 & 5.0\\
Near boundary & Yes & 96.7 & 3.3\\
Hidden slow memory & No & 98.3 & 11.7\\
Correlated readings & No & 98.3 & 1.7\\
Biased readings & No & 0.0 & 100.0\\
Behaviour change & No & 95.0 & 5.0\\
\bottomrule
\end{tabular}
\end{center}
The nominal one-decision band coverage is 95 percent under the valid contract;
these small, dependent-budget frequencies are diagnostics, not a new coverage
theorem. A tap-band hit concerns the observed response, not membership in the
finite candidate class or stability of future behaviour. In particular, hidden
slow memory falls outside the declared class even when its first 12 taps fit
the bands. Under behaviour change the readings describe the old response.
Biased readings empty the set throughout this run and force the certificate
policy to abstain. Correlated readings use an unverified standard-error model.
None of these four regimes retains the stated guarantee; a high observed band
hit or zero observed financial breaches does not make an invalid contract valid.

\subsection*{Pending evidence: deliberately empty result cells}
The following numerical cells remain empty. The returned archive contained no
market CSV, dynamic-inventory experiment, institution-level intervention or
external replication. The general theorem and independent novelty review also
remain open; a table placeholder is not a theorem statement or a proof.

\begin{center}\small
\begin{tabular}{p{0.62\textwidth}cc}
\toprule
Pending study & Estimate & Uncertainty\\
\midrule
Untouched market-data test beyond order-flow imbalance and depth & & \\[1.8ex]
Full changing-inventory, execution and funding comparison & & \\[1.8ex]
Dealer/venue intervention with endogenous market response & & \\[1.8ex]
Net value under empirically calibrated disclosure costs & & \\[1.8ex]
Independent implementation and replication & & \\[1.8ex]
\bottomrule
\end{tabular}
\end{center}
\noindent Further data cannot by themselves prove a general identification or
decision-value theorem. Its exact statement, proof and literature priority
require separate review. No publication rating follows from the returned run.

% END INFORMATION DECISION REVISION 20261029

% BEGIN FINITE HORIZON COST 20261029
\section*{Finite-horizon information, cohort ambiguity, and net inventory-service value}

\paragraph{Status and contribution boundary.}
This appendix supplies a restricted analytic separation and its financial
consequence. The proof below is written in full; exact rational and numerical
checks accompany it. It has not received independent proof review and is not
a Lean formalization. Its priority is unconfirmed. Control without unique
identification, task-directed exploration, Gaussian functional estimation,
moment matching, and the value of information are established ideas
\citep{fhc:waarde,fhc:wagenmaker,fhc:donoho,fhc:heinrich,fhc:blackwell}.
The candidate extension is their explicit combination for a finite-horizon
inventory commitment, with a hard stress budget and the information bill
deducted from both capital and profit. It does not close the broader open
theorem of the preceding drafts.

\subsection*{The observation contract and the financial institution}
Fix $H\ge4$, $0\le a<1$, $\kappa>0$ and
$0\le Q\le\overline Q<(1-a)/\kappa$. There are arbitrarily many finite
cohorts, with $q_i\ge0$, $\sum_iq_i=Q$, and $0<\eta_i\le1$:
\[
 x_{t+1}=a x_t+\kappa\bigl(\sum_iq_i y_{i,t}+u_t\bigr),\qquad
 y_{i,t+1}=(1-\eta_i)y_{i,t}+\eta_i x_{t+1}.
\]
Initial states are reset to zero. Let $h_0,\ldots,h_{H-1}$ be the price
response to a unit impulse, with $h_0=\kappa$, and put
$G=\sum_{t<H}h_t$. Positivity implies $h_t\ge0$.
No positive lower bound on $\eta_i$ is imposed.

The observer purchases $N$ independent reset experiments, each lasting $H$
model periods and returning all $H$ taps with independent Gaussian errors of
known variance $\sigma^2$. Thus the tap means are
$\overline h_t=h_t+\sigma Z_t/\sqrt N$, with independent standard normal
$Z_t$. The same observations are available for identifying the cohorts and
for choosing the financial position. One experiment is not one scalar reading:
it returns $H$ readings and costs $c_H$, including the declared operating and
time charges. Total information cost is $F=Nc_H$. This observation contract
is not an ordinary exchange tape; longer horizons cost both time and readings.
Plant parameters, the feedback cap, reset, and stationarity are assumptions.
Measurements are acquired ahead of a common future service-start deadline.
Cash can earn the common reference rate during acquisition; the model does
not include earlier service opportunities. The stated measurement price is
exogenous, not a calibrated self-impact cost for actual market probes.

The institution is a price-taking, unlevered inventory-warehousing desk.
It buys $I\ge0$ units at $P_0$, holds them for $H$ periods, and closes.
It earns a prequoted service fee $pI$ \emph{after} settlement. Round-trip
execution and carry costs are $bI+dI^2$, where $b\ge0$ and $d>0$.
For example, $b=b_{\rm trade}+P_0 r_f H/(365\cdot24\cdot60)$ for
one-minute model periods. The service premium is unavailable as collateral
until after settlement. Both trading costs and $F$ must be prefunded.
Initial cash is $W$, the stressed-settlement loss budget is $B$, and
$|u_t|\le\overline u$, with $\overline u>0$. Assume
$P_0>\overline u\,\overline G$, so the asset price remains positive.
This desk does not alter market feedback, depth, or another participant's
behaviour. The service fee is exogenous, not an equilibrium price.

For the expected-profit objective, future shocks have conditional mean zero
given the observations. Consequently the conditional expected trading price
change is zero. The financial objective, including every declared charge, is
\[
 V(I,F)=g(I)-F,\qquad g(I)=(p-b)I-dI^2.
\]
This is expected service profit with a hard stress constraint, not a Kelly
growth theorem or a return-prediction result. There are no additional training
or validation tapes, no leverage, uncertain fills, or changing inventory.

\begin{lemma}[Sharp prior bounds and exact prefunded capacity]\label{fhc:capacity}
Write
\[
 \underline G=\kappa\sum_{t<H}a^t,\qquad
 \overline G=\kappa\sum_{t<H}(a+\kappa\overline Q)^t.
\]
Then $\underline G\le G\le\overline G$, and both endpoints are attained
in the declared class. For a gain bound $C\ge G$ and
$0\le F\le\min(B,W)$, define
\[
 \begin{gathered}
 R(v,r)=\frac{2v}{r+\sqrt{r^2+4dv}},\qquad I_u=\frac{(p-b)_+}{2d},\\
 I(C,F)=\min\{I_u,R(B-F,\overline u C+b),R(W-F,P_0+b)\}.
 \end{gathered}
\]
This is the optimal expected-profit position subject to the stated stress
reserve and prefunding constraints. The zero-violation no-observation robust optimum is
$I_{\rm nd}=I(\overline G,0)$, not an arbitrarily conservative benchmark.
The response-informed oracle is $I_\circ=I(G,0)$; it knows no future shocks.
\end{lemma}

\begin{proof}
Let $k_j=\sum_iq_i\eta_i(1-\eta_i)^j$. The normalized price response is
the renewal sequence with defective interarrival probabilities
$s_1=a+\kappa k_0$ and $s_j=\kappa k_{j-1}$ for $j\ge2$.
Their total mass is $a+\kappa Q<1$, with the remaining mass placed at
infinity. The cumulative interarrival distribution at integer $j\ge1$ is
\[
 a+\kappa\sum_iq_i[1-(1-\eta_i)^j].
\]
It lies between $a$ and $a+\kappa\overline Q$. Coupling interarrivals by
the same independent uniforms orders every renewal time and hence the
number of renewals before $H$. Taking expectations proves the two bounds.
The lower endpoint is the no-feedback economy; the upper is one cohort
with $Q=\overline Q$ and $\eta=1$. This also proves sharpness of the
no-observation stress bound.

At date $t\le H$, the price change is
$x_t=\sum_{j<t}h_j u_{t-1-j}$, so
$|x_t|\le\overline u\sum_{j<t}h_j\le\overline uG$.
All costs paid by that date are at most the total prefunded costs.
The constant adverse sequence $u_t=-\overline u$ attains the total price
loss $I\overline uG$ at settlement. Before the service fee is credited,
worst loss from initial wealth is therefore exactly
$I\overline uG+bI+dI^2+F$. The declared robust reserve requirement is
$I\overline uC+bI+dI^2+F\le B$; prefunding requires
$P_0I+bI+dI^2+F\le W$. These increasing quadratics yield the two roots.
Finally $g$ increases up to $I_u$ and decreases after it, proving the
optimizer and the two benchmarks. This exact reserve contract is not a
peak-to-trough drawdown or a repeated-deployment guarantee.
\end{proof}

\begin{lemma}[No data at the same error budget]\label{fhc:nodata-chance}
Let $I_f=\min\{I_u,R(W,P_0+b)\}$ and allow a data-free rule to randomize,
with probability of a stress-budget violation at most $\gamma$ in every
economy of the class. Its greatest expected service profit is
\[
 V_{{\rm nd},\gamma}=(1-\gamma)g(I_{\rm nd})+\gamma g(I_f).
\]
It is attained by the displayed two-position lottery. This is a theoretical
benchmark, not a recommended trading policy. The corresponding response-
informed same-error-budget oracle has value
$(1-\gamma)g(I_\circ)+\gamma g(I_f)$.
\end{lemma}

\begin{proof}
Truncating an action at $I_f$ weakly improves profit and lowers risk, so
restrict to $0\le I\le I_f$. At the attainable worst economy with gain
$\overline G$, every $I>I_{\rm nd}$ is unsafe unless $I_{\rm nd}=I_f$,
in which case all larger actions are already dominated. Thus a feasible
lottery puts at most $\gamma$ mass above $I_{\rm nd}$. Below that point
profit is at most $g(I_{\rm nd})$, and above it at most $g(I_f)$.
This gives the upper bound. The two-point lottery attains it and is unsafe
in any economy with probability at most $\gamma$. Replace $\overline G$
by the known $G$ for the oracle statement.
\end{proof}

\begin{theorem}[Uniform cost-aware finite-horizon decision]\label{fhc:decision}
Fix $0<\gamma<1/2$, $z=\Phi^{-1}(1-\gamma)$,
$s_N=\sigma\sqrt{H/N}$, and
\[
 Y=\sum_{t<H}\overline h_t,\qquad
 C_N=\min\{\overline G,\max\{\underline G,Y+zs_N\}\}.
\]
Purchase a fixed, affordable $N$ and take $I_N=I(C_N,F)$.
The following hold uniformly over the declared positive-cohort class,
without dependence on cohort spacing, cohort count, or a memory-rate floor.
\begin{enumerate}
\item With probability at least $1-\gamma$ over measurements, the position
      satisfies the stress budget for \emph{every} bounded future shock path.
\item Put $r_0=\overline u\underline G+b$, $I_{\max}=W/P_0$,
\[
 \begin{aligned}
 L_G&=\frac{\overline u I_{\max}}{r_0}, &
 L_F&=\frac{1}{\min(r_0,P_0+b)},\\
 A&=2z\sigma\sqrt H\,(p-b)_+L_G, & D&=1+(p-b)_+L_F.
 \end{aligned}
\]
With probability at least $1-2\gamma$,
\[
 0\le g(I_\circ)-V(I_N,F)\le A N^{-1/2}+D c_HN.
\]
For sufficiently small $c_H>0$, choosing
$N=\lceil(A/(2Dc_H))^{2/3}\rceil$ is affordable and gives a bound
$3\,2^{-2/3}A^{2/3}(Dc_H)^{1/3}+Dc_H$, when $A>0$.
This is an upper bound from this procedure, not a new minimax optimality claim.
\item Let $I_-=I(G+2zs_N,F)$. A sufficient condition for strictly positive
      expected net information value relative to the strongest no-data rule
      at the \emph{same} error budget is
\[
 \underbrace{(1-2\gamma)g(I_-)-F-V_{{\rm nd},\gamma}}_{\displaystyle\mathcal L_N}>0.
\]
Precisely, $\mathbb E_D[V(I_N,F)]-V_{{\rm nd},\gamma}\ge\mathcal L_N$.
The conclusion can fail at high data costs or a low service premium.
\end{enumerate}
\end{theorem}

\begin{proof}
$Y-G$ is Gaussian with variance $H\sigma^2/N$. On
$\{Y-G\ge-zs_N\}$, which has probability $1-\gamma$, $C_N\ge G$;
clipping preserves this because $G$ lies between the prior bounds.
Lemma~\ref{fhc:capacity} proves simultaneous safety for every future
bounded path on this single event, including the data-selected position.
There is no tail beyond $H$ to estimate, explaining the absence of a
memory-floor assumption. This does not certify an infinite horizon.

On $E=\{|Y-G|\le zs_N\}$, of probability $1-2\gamma$,
$G\le C_N\le G+2zs_N$. The position function is decreasing in each
of $C$ and $F$, so $I_-\le I_N\le I_\circ$ on $E$.
To bound the position difference, implicitly differentiate the stress root:
its derivatives are $-\overline u I/(\overline u C+b+2dI)$ in $C$
and $-1/(\overline u C+b+2dI)$ in $F$.
Clip this root at $I_{\max}$ without changing the final minimum, since
the prefunding root is at most $I_{\max}$. The clipped root is
$L_G$-Lipschitz in $C$ and $1/r_0$-Lipschitz in $F$.
The prefunding root is independent of $C$ and
$1/(P_0+b)$-Lipschitz in $F$. Minimum with $I_u$ preserves these
common Lipschitz bounds. Hence
\[
 0\le I_\circ-I_N\le 2L_Gzs_N+L_FF.
\]
On $[0,I_u]$, $0\le g'\le(p-b)_+$. Multiplication and addition of
the paid fee give the regret bound. Minimizing
$AN^{-1/2}+Dc_HN$ over positive real $N$ gives the stated root and
value. Rounding up adds at most $Dc_H$; affordability follows as
$c_HN=O(c_H^{1/3})\to0$. The case $A=0$ has no positive-gain
information task and does not need this optimization.

On $E$, $g(I_N)\ge g(I_-)$, while $g(I_N)\ge0$ always because
$0\le I_N\le I_u$. Thus
$\mathbb E_D g(I_N)\ge(1-2\gamma)g(I_-)$.
Subtracting $F$ and the same-error-budget no-data value proves the last assertion.
Conditional mean-zero future flows justify the stated expected-profit
interpretation; their distribution is not learned from extra tapes.
\end{proof}

\begin{theorem}[A two-cohort identification obstruction under the same observations]\label{fhc:pair}
Fix $Q>0$, $a+\kappa Q<1$, and $0<\eta_0<1$. For sufficiently small
$\delta>0$, compare the two economies
\[
\begin{array}{c|cc}
 &\text{normalized weights}&\text{adaptation rates}\\ \hline
\mathrm A_\delta &(1/2,1/2)&(\eta_0-\delta,\eta_0+\delta)\\
\mathrm B_\delta &(1/5,4/5)&(\eta_0-2\delta,\eta_0+\delta/2).
\end{array}
\]
Actual weights are $Q$ times those displayed. For fixed $H\ge4$,
$D_H(\delta)=\|h^{\rm B}_{0:H-1}-h^{\rm A}_{0:H-1}\|_2
=\Theta(\delta^3)$. The exact minimax error for distinguishing these
two \emph{known} alternatives from $N$ reset experiments is
\[
 e_N=\Phi\!\left(-\frac{\sqrt N D_H(\delta)}{2\sigma}\right).
\]
Therefore error at most $\beta<1/2$ requires
$N\ge[2\sigma\Phi^{-1}(1-\beta)/D_H(\delta)]^2
=\Theta(\delta^{-6})$.
In particular any recovery of the two normalized weights to sup-norm error
less than $0.15$ with probability at least $1-\beta$ in both worlds must
satisfy this testing requirement.
Both worlds contain exactly two distinct rates. Both have identical infinite
summed response, but their finite-horizon gains are not asserted identical.
\end{theorem}

\begin{proof}
The two rate distributions have equal mass, mean $\eta_0$, and variance
$\delta^2$. Their third centered moments are respectively $0$ and
$-3\delta^3/2$. Thus their kernels $k_j$ agree at $j=0,1$ and
$k^{\rm B}_2-k^{\rm A}_2=-3Q\delta^3/2$. The response recursion
gives equal $h_0,h_1,h_2$ and
\[
 h^{\rm B}_3-h^{\rm A}_3=-3\kappa^2Q\delta^3/2,
\]
which is a lower bound for $D_H$ in absolute value.

For an upper bound put $f_t(\eta)=\eta(1-\eta)^t$.
Its third derivative is
$3t(t-1)(1-\eta)^{t-2}
-\eta t(t-1)(t-2)(1-\eta)^{t-3}$, with absent low-degree terms zero.
The sum over $t\ge0$ of its absolute value is at most $12/\eta^3$.
For $\delta\le\eta_0/4$, all intervening rates are at least
$v=\eta_0/2$. Taylor's integral remainder in $\ell_1$ is therefore at
most $2|\eta-\eta_0|^3/v^3$. The matching first two moments cancel
the quadratic Taylor terms. The sum of the two third absolute centered
moments is $(1+17/10)\delta^3=27\delta^3/10$, giving
\[
 \|k^{\rm B}-k^{\rm A}\|_1\le\frac{27Q}{5v^3}\delta^3.
\]
If $K(z)=\sum k_tz^t$ and $H(z)=\sum h_tz^t$, the recursion gives
$H(z)=\kappa/[1-az-\kappa zK(z)]$. Consequently
$H_{\rm B}-H_{\rm A}=zH_{\rm A}H_{\rm B}(K_{\rm B}-K_{\rm A})$.
The infinite response sums of both economies equal
$S=\kappa/(1-a-\kappa Q)$. Young's inequality yields
$D_H\le S^2\|k^{\rm B}-k^{\rm A}\|_1=O(\delta^3)$.
This proves the two-sided order with constants independent of $\delta$.

For equal-covariance Gaussian laws, the likelihood-ratio test projects the
sample mean onto the difference of the two known means, and thresholds at
their midpoint. The projected normal law gives the displayed error in each
world. Its equal-prior Bayes optimality and symmetry establish minimax
optimality. A weight estimator accurate within $0.15$ distinguishes the
two disjoint normalized-weight balls, proving the identification lower bound.
Finite-horizon gains need not be equal: only their difference is of cubic
order. The test concerns latent cohort recovery, not identification of the
input-output response vector, and does not cover alternative probe designs.
\end{proof}

\begin{corollary}[A priced separation]\label{fhc:priced}
Fix $c_H>0$ and an affordable $N$ for which $\mathcal L_N>0$ at the limiting merged
economy $\delta=0$. For all sufficiently small positive $\delta$, the
financial rule of Theorem~\ref{fhc:decision} has a
strictly positive expected net information value, with its stated safety
probability. Its observation requirement has no dependence on $\delta$.
At this same fixed information price, sufficiently small $\delta$ makes cohort
identification at error $\beta$ unaffordable within the same hard fee budget
$B$: $c_H[2\sigma\Phi^{-1}(1-\beta)/D_H(\delta)]^2>B$.
Any procedure limited to that many affordable observations is covered,
including one that stops early, since it can be viewed as a function of the
full affordable sample. A rational agent may abandon identification and use
the no-data rule. There is no assertion that regulation requires identification.
\end{corollary}

\begin{proof}
Apply the two theorems. The first statement uses the strictly positive
lower bound on net value, rather than an assumption that more information
is automatically beneficial. The second follows from $D_H=\Theta(\delta^3)$.
For a neighbourhood of a fixed $\delta$ where the strict value inequality
holds, continuity preserves it. To claim the inequality for all sufficiently
small $\delta$, it suffices that it hold strictly at the limiting merged
economy: finite $G_H$ and $I_-$ are continuous there. The numerical example
reports both the primary pair and that limiting check. No finite-class grid
is used in the uniform decision guarantee.
\end{proof}

\paragraph{What the result does not establish.}
The direct gain-sum certificate is exactly ordinary Gaussian projection onto
the vector of ones. It is not a new estimator and does not beat that benchmark.
The latent-cohort lower bound is not a lower bound for learning the observable
response vector; comparing it with FIR model-identification results for
$\mathcal H_\infty$ norms would conflate different tasks \citep{fhc:tu}.
No optimal estimation rate over all observation designs is claimed.
The data cost is not an empirically measured exchange probing cost.
An ordinary public-tape bridge, endogenous quotes and depth, incomplete resets,
unknown or misspecified noise, and independent replication remain open.
The information purchase in the example is predeclared. Its positive value
is pointwise at the displayed economy; no prior over unknown economies or
globally optimal Bayesian/robust acquisition policy is claimed.
The fixed horizon, independent measurement contract, and conditional mean-zero
future flows cannot be removed by the simulation. At another horizon the
certificate must be recomputed. A rolling sequence needs simultaneous or
time-uniform coverage; a one-decision guarantee is insufficient.

\paragraph{Literature priority.}
The reviewed sources already establish the general principles used here.
The explicit dollar-cost/capital formulation and the combined result above
are candidate additions to this project, not verified first-in-literature
claims. A specialist must check whether they are an immediate corollary of
existing task-specific or constrained-decision results before publication.

\begin{thebibliography}{9}\small
\bibitem[van Waarde et~al.(2020)]{fhc:waarde} H.~J. van Waarde, J. Eising,
H.~L. Trentelman and M.~K. Camlibel. Data Informativity: A New Perspective on
Data-Driven Analysis and Control. \emph{IEEE TAC} 65, 4753--4768.
\bibitem[Wagenmaker et~al.(2021)]{fhc:wagenmaker} A.~J. Wagenmaker,
M. Simchowitz and K. Jamieson. Task-Optimal Exploration in Linear Dynamical
Systems. \emph{ICML}, PMLR 139, 10641--10652.
\bibitem[Donoho(1994)]{fhc:donoho} D.~L. Donoho. Statistical Estimation and
Optimal Recovery. \emph{Annals of Statistics} 22, 238--270.
\bibitem[Heinrich and Kahn(2018)]{fhc:heinrich} P. Heinrich and J. Kahn.
Strong Identifiability and Optimal Minimax Rates for Finite Mixture Estimation.
\emph{Annals of Statistics} 46, 2844--2870.
\bibitem[Blackwell(1953)]{fhc:blackwell} D. Blackwell. Equivalent Comparisons
of Experiments. \emph{Annals of Mathematical Statistics} 24, 265--272.
\bibitem[Tu et~al.(2018)]{fhc:tu} S. Tu, R. Boczar and B. Recht.
Minimax Lower Bounds for $\mathcal H_\infty$-Norm Estimation.
\emph{arXiv:1809.10855}.
\end{thebibliography}

\section*{Checked synthetic financial illustration}
This illustration uses the frozen primary design and the stronger comparison
recorded in the design amendment. It is exploratory and was developed after
the earlier project runs; it is not independent preregistration or replication.
The two-cohort primary has $a=0.4$, $\kappa=1$, $Q=0.4$, rates
$0.49,0.51$, and equal weights. The known class cap is $\overline Q=0.55$;
$H=60$ one-minute model periods, $N=32$, $\sigma=0.02$, and
$\gamma=0.025$. Initial wealth is 1,000, price 100, stress budget 1,
and flow bound 0.04. The service fee is 0.08 per unit. Execution costs are
$0.002I+0.0002I^2$ and carry is charged at 5 percent annualized.
The assumed all-in reset fee is 0.002, so the learning bill is 0.064.
These are hypothetical model dollars, not an exchange calibration or returns.
The same-error-budget no-data lottery is a theoretical strongest comparator,
not a proposed desk practice. The zero-violation comparisons are also shown.

\begin{center}\footnotesize
\begin{tabular}{p{0.43\textwidth}rrr}
\toprule
Policy & Mean units & Net expected profit & Unsafe reserve (\%)\\
\midrule
No data, zero violations & 1.306 & 0.1007 & 0.00\\
Best no data, same error budget & 1.523 & 0.1171 & 2.50\\
Finite-horizon certificate & 4.553 & 0.2844 & 2.44\\
Standard Gaussian projection & 4.553 & 0.2844 & 2.44\\
Full-response ellipsoid & 4.386 & 0.2718 & 0.00\\
Point estimate (uncertified) & 4.601 & 0.2880 & 50.41\\
Exact gain report, fee 0.05 & 4.669 & 0.3072 & 0.00\\
Response oracle, zero violations & 4.914 & 0.3756 & 0.00\\
Response oracle, same error budget & 5.041 & 0.3851 & 2.50\\
\bottomrule\end{tabular}\end{center}
Profit integrates the future zero-mean flow model and every declared cost;
the data-dependent policy values are also integrated over reading noise by
one-dimensional Gaussian quadrature. Mean units and empirical reserve
violations use 20,000 synthetic reading draws. Oracle and no-data lottery
values are analytic. A reserve violation means some bounded future path can
exceed the budget, not that a sampled market path actually did so.
The 97.5 percent safety statement is over measurement draws for a single
future commitment; it is not certainty, market-wide stability, or time-uniform
coverage. The point rule buys slightly more expected profit by violating the
same risk requirement and is not a valid superiority comparison.
The proposed estimator and standard Gaussian projection coincide exactly.

The net information advantage over the strongest same-error-budget no-data rule is 0.1673. The written theorem's sufficient lower bound is 0.1465, strictly positive even after all costs. At the limiting merged economy the lower bound is 0.1465; continuity therefore supports the small-separation corollary. These bounds are evaluated in floating point with a wide positive margin, not claimed interval-arithmetic certificates.
\begin{center}\small\begin{tabular}{rrr}
\toprule
Separation $\delta$ & Known-pair resets required & Implied fee bill\\
\midrule
0.0800 & 478 & 0.955\\
0.0400 & 3.96e+04 & 79.2\\
0.0200 & 2.84e+06 & 5.67e+03\\
0.0100 & 1.91e+08 & 3.83e+05\\
0.0050 & 1.26e+10 & 2.51e+07\\
0.0025 & 8.15e+11 & 1.63e+09\\
\bottomrule\end{tabular}\end{center}
The known-pair test receives all 60 taps, exactly as the financial rule does.
Response differences are calculated in rational arithmetic before conversion
to floating point. The sample threshold is the Gaussian testing formula, not
a measured runtime or an asserted physical trading-capital requirement.
For the primary $\delta=0.01$, distinguishing the two already-known
alternatives requires about 191 million resets at 5 percent error. Even
this optimistic information fee is far beyond the unit stress budget.
Identifying arbitrary cohorts cannot be easier than this embedded test under
the same fixed query contract. Other experimental designs are not covered.

At fixed $N=32$, the expected net benefit crosses zero at an illustrative reset price of about 0.005825. An exact scalar gain report is worth purchasing only below a fee of about 0.188690 in this example. These are computed break-even prices, not observed reporting costs.
\paragraph{Failures retained.}
The full output includes all 819 cost, horizon and service-fee sensitivity
rows, including negative values and unaffordable information budgets.
A low service premium or expensive measurements can make abstaining from
data acquisition preferable. Independent checks cover 13 algebra, optimizer,
shock-corner, cash-flow and Gaussian-testing properties, including the
same-error-budget no-data optimum checked against a separate linear program.
An initial low-premium LP check was affected by the solver's absolute
objective tolerance; rescaling the objective resolved it, and that failed
log is retained. This is not independent mathematical proof review.

\begin{center}\footnotesize\begin{tabular}{p{0.52\textwidth}rr}
\toprule
Reading/model contract & Gain coverage (\%) & Unsafe reserve (\%)\\
\midrule
Independent Gaussian, valid & 97.36 & 2.63\\
Correlated taps, wrong variance & 71.16 & 28.84\\
Correlated taps, correct variance & 97.36 & 2.63\\
Negative reading bias, invalid & 0.00 & 100.00\\
Future feedback changes, invalid & 0.00 & 100.00\\
\bottomrule\end{tabular}\end{center}
The valid Gaussian row has finite Monte Carlo variation around its nominal
97.5 percent coverage. Using the correct correlated variance restores the
specified measurement contract at the cost of a larger margin. Bias or
future parameter change is not repaired by a larger sample under the wrong
model. A signed kernel such as $(1,-0.9)$ has sum $0.1$ and absolute sum
$1.9$, so positivity is indispensable to using the plain sum. The same
certificate used repeatedly without correction does not retain its
one-decision guarantee.

\paragraph{What changed economically.}
In this model, the desk can choose a larger inventory commitment while
paying for information and retaining the same permitted unsafe-action
probability as its strongest no-data comparator. This is an example of
positive decision value, not evidence that information always pays, that
trading strategies must be identified, or that a real venue's liquidity and
welfare improve. Realistic observation costs and a public-data bridge remain
requirements for an empirical finance claim.

% END FINITE HORIZON COST 20261029


% BEGIN ADVANCED RETURN STATUS 20261030
\clearpage
\section*{Returned advanced experiment: evidence and limits}
\addcontentsline{toc}{section}{Returned advanced experiment: evidence and limits}
\begingroup
\renewcommand{\baselinestretch}{1.08}\selectfont
\paragraph{Research question and financial institution.}
Can information about an aggregate market response improve an exposure decision
after information, execution and carry costs, while the hidden trading cohorts
remain difficult to recover? The new experiment concerns a price-taking,
unlevered inventory-warehousing desk. It buys a fixed inventory, holds it for
$H=10,60,240$ one-minute \emph{model} periods and closes. A prequoted service
premium arrives after settlement. It cannot change market feedback, quotes or
other traders. This is a separate application from a dealer or venue that
changes admission and thereby changes the market itself. The objective is
expected service profit under a stress reserve and cash constraint, not
logarithmic portfolio growth or predicted crypto returns.

\paragraph{Completed execution record.}
The author returned the Deep Research archive with base seed 20261101.
All 611 manifest entries, all 288 checkpoint-to-statistics links and all 15
released source files were verified. The 288 generated economies use 24 seed
clusters, 4,096 measurement repetitions per economy and $N$, $N=8,32,128$,
three information-price settings, seven measurement rules and a same-error
oracle. The 20,736 policy/cost rows share readings and are not independent
market experiments. All 27 implemented tests passed. Numeric policy fields
reproduce the local same-seed run to about $1.25\times10^{-14}$. This is
cross-environment reproducibility of the same implementation, not independent
replication. The returned runtime used Python 3.13.18 and the pinned numerical
packages; the research runner took 52.71 seconds, excluding notebook setup.

\paragraph{What Lean checked.}
Lean 4.19.0 checked seven concrete rational inequality certificates, with no
reported axioms. They include a conservative positive-value numerical witness
greater than 0.14 hypothetical model dollars for the original primary contract.
They do not formalize the Gaussian analytic tail step, market dynamics, general
separation theorem, translation from observations to economics or novelty.
The earlier positive illustration is retained as a pointwise example. It must
not be read as a positive average-value theorem over all economies or prices.

\paragraph{Economic value depends on the information bill.}
For the standard Gaussian rule within the independent-Gaussian economy design:
\begin{center}\small
\setlength{\tabcolsep}{4pt}
\begin{tabular}{@{}p{0.28\linewidth}rrr@{}}
\toprule
Information price & Mean advantage & 95\% interval & Positive / total\\
\midrule
Free & $+0.13047$ & $[0.08458,0.18420]$ & 99 / 144\\
Baseline & $-0.02953$ & $[-0.07022,0.01310]$ & 59 / 144\\
Five times baseline & $-0.17212$ & $[-0.21856,-0.12715]$ & 26 / 144\\
\bottomrule
\end{tabular}
\end{center}
The comparator is the analytically optimal no-data rule with the same permitted
unsafe-action probability, $\gamma=0.025$. Values are hypothetical dollars per
inventory-service decision. Baseline cost is
$N(0.0005+0.000025H)$; its dependence on horizon and $N$ is explicit, not an
exchange calibration. The intervals use 5,000 percentile bootstrap resamples
of 24 seed-cluster averages. This price-specific analysis was added after the
return and is exploratory, not preregistered confirmation or uncertainty about
actual market profit. The baseline-price interval crosses zero: this design
does not establish a positive average benefit after costs. The high-price
interval is negative under the stated simulator sampling scheme.

\paragraph{Protection and profitability are different outcomes.}
For purchased iid-Gaussian rules in iid-Gaussian economies, the mean simulated
unsafe-action frequency was 0.834\%, below the 2.5\% allowance. Across 6,535
declared-valid purchased-rule checks, Holm correction produced zero rejections.
Individual finite-sample frequencies can exceed the allowance; non-rejection
is not a proof. An unsafe action means its reserve can breach the budget under
a permitted worst-case shock path. It is not an observed liquidation, default
or historical crash. The Gaussian and declared-variance rules coincide exactly
under iid readings. Their advantage over a full-response ellipsoid is a
task-specific comparison of established methods, not a new estimator.

\paragraph{Failures retained and interpreted.}
\begin{center}\small
\setlength{\tabcolsep}{4pt}
\begin{tabular}{@{}p{0.66\linewidth}r@{}}
\toprule
Purchased Gaussian rule: observation/future setting & Mean unsafe action\\
\midrule
Independent Gaussian readings, stationary future & 0.834\%\\
Correlated readings, iid variance used & 14.695\%\\
Downward-biased readings & 31.169\%\\
Future feedback changes & 61.033\%\\
\bottomrule
\end{tabular}
\end{center}
These are unweighted descriptive averages of dependent configuration rows,
including the zero-price setting. They are not pooled market failure
probabilities. Under correlated Gaussian readings, a declared-variance rule
reduces the mean to 0.409\%, and a projected Student-$t$ rule to 1.153\%.
Protection can improve while the information purchase remains uneconomic.
The Student-$t$ guarantee requires independent Gaussian reset totals; good
observed performance on $t_3$ noise does not extend that guarantee. Median of
means needs independent resets and an externally valid variance bound. At
insufficient $N$ it uses a paid prior-bound fallback. Neither robust interval
repairs mean bias or a change in future dynamics. The full design retains
7,489 negative-advantage rows and 1,743 unaffordable purchase rows. Pooling all
methods and scenarios into one success percentage would be misleading.

\paragraph{Information acquisition is an open economic decision.}
The fixed-$N$ protocol buys information when affordable; it does not optimize
the purchase. In 32 of 288 generated economies the known service premium is
no greater than known linear execution/carry cost. Positive inventory then has
no positive expected service surplus before any information is obtained.
In the first indexed case, s000\_c000 at $N=32$ and baseline price, no-data and
response-informed oracle profit both equal approximately 0.75902, because
response information cannot relax the relevant position limit. Paying 0.024
for information leaves approximately 0.73500. A useful measurement need not
change the financial action. A purchase rule should use only information
available before purchase, include the option to abstain, and avoid selecting
$N$ or a policy after looking at hidden gains or test outcomes. That rule is
still to be proved and validated; it has not been added to this returned run.

\paragraph{Observation boundary and contribution status.}
There are no historical coins, exchange tapes or crash replays in this run.
The quote-data stage is \textbf{NOT\_RUN}; the unit-test quote fixture is
synthetic. Controlled resets, known plant/noise inputs, a feedback cap and
stationarity are maintained assumptions. The finite-horizon positive-response
certificate needs no adaptation-rate floor when every relevant tap is observed
under that reset contract. This does not remove the hidden-tail problem for
infinite-horizon or ordinary public-tape certification. New observations do not
independently verify the earlier passive, capital, AI or equilibrium claims.

The supported interpretation is a conditional financial separation: aggregate
response information can support a specified exposure decision without latent
cohort recovery, and positive net value requires additional economic conditions.
The written general proof still needs independent review and its distinction
from existing identification, data-informativity and decision theory needs
specialist assessment. No publication-quality rating follows from execution.

\paragraph{Traceability.}
The immutable experiment source is repository commit \texttt{6a319b0db147}.
The return audit is in commit \texttt{c0c0c010ec00}. The study directory and
its relative audit directory are shown separately for readability:
\begin{quote}\footnotesize
\path{studies/advanced_cost_separation_20261029/}\\
\path{validation/returned_colab_20261030/}
\end{quote}
The uploaded archive is:
\begin{quote}\footnotesize
\path{MM_advanced_deep_seed20261101_v1.zip}
\end{quote}
Its SHA256 begins \texttt{e35c10f6b276}; the full hash, configuration, source
fingerprint and exploratory price/paired audits are retained in that directory.
No original return outcome has been overwritten.
\endgroup

\clearpage
\section*{Research work register: 30 October 2026}
\addcontentsline{toc}{section}{Research work register: 30 October 2026}
\begingroup
\renewcommand{\baselinestretch}{1.08}\selectfont
This register is a dated snapshot, not a claim that the project is complete.
\emph{Completed} means the named artifact or check exists. \emph{Written,
review open} means an argument is present but has not received independent
mathematical sign-off. \emph{Pending} means the evidence or proof is still
required. Good outcomes from another experiment do not silently close a task.

\subsection*{Completed records to preserve}
\begin{center}\small
\setlength{\tabcolsep}{4pt}
\begin{tabular}{@{}p{0.13\linewidth}p{0.73\linewidth}@{}}
\toprule
Record & Verified completion and boundary\\
\midrule
C1 & Returned archive, source fingerprint and 288 checkpoint links verified. Same-code reproducibility; no independent implementation.\\
C2 & All 27 implementation checks passed. They do not test every theorem in the papers.\\
C3 & Seven numerical Lean certificates passed. The full general theorem is not formalized.\\
C4 & 288-economy synthetic audit completed; every negative and unaffordable row retained. No market-data stage ran.\\
C5 & Pointwise positive-value illustration retained; broad price-sensitive results and assumption failures recorded. No universal profitability claim.\\
C6 & Targeted literature comparison and a written restricted separation exist. Originality and complete proof review remain open.\\
\bottomrule
\end{tabular}
\end{center}

\subsection*{Core work before stronger contribution claims}
\paragraph{R1. Establish the precise literature gap. Status: pending specialist review.}
Compare the exact assumptions, observation experiment, decision target and
conclusion with identification theory, data informativity, task-optimal
exploration, functional estimation and constrained financial decisions. Separate
latent cohort recovery from learning the observable response. Do not claim
that control or allocation without complete identification is new.
\emph{Completion evidence:} a source-backed comparison table identifying a
specific non-overlap, or a clear statement that the result is an application
or corollary; resolved reviewer comments. More simulations cannot complete R1.

\paragraph{R2. Independently audit the load-bearing proofs. Status: written, review open.}
Review the finite-cohort renewal bounds, observation-to-confidence implication,
uniformity in hidden memory, fixed-query testing obstruction, all fee/cash
constraints and the same-error no-data optimum. Audit the older broad
certification, passive-identification and capital claims separately. Check
quantifiers, constants, admissible probes, reset assumptions and counterexamples.
\emph{Completion evidence:} theorem-by-theorem assumptions and proof register,
resolved objections and independent mathematical review. Record each Lean
theorem's actual scope; do not substitute seven scalar facts for a general proof.

\paragraph{R3. Make the information purchase economically explicit. Status: pending.}
Define what the desk knows before buying observations and whether the acquisition
objective is robust, Bayesian or another stated criterion. Include abstention
and screen cases where the premium cannot cover known costs or information
cannot change the constrained position. Select sample size and price using
available information, not hidden gains or test outcomes.
\emph{Completion evidence:} a precise acquisition rule, a valid decision-value
or regret statement, matched-information/cost comparisons, and a frozen new
validation design. Reusing this return to design a rule is exploratory.

\paragraph{R4. Justify the financial institution and measurement contract. Status: partial.}
The fixed inventory, horizon, post-settlement premium and accounting are explicit.
The economic source of the premium, operating cost of reset observations,
information reliability and acquisition opportunity cost need justification.
An own-account investor, warehousing desk, market maker and venue supervisor
have different actions and incentives. Keep their claims separate.
\emph{Completion evidence:} a coherent institutional mechanism and documented
parameter/cost interpretation, with comparative statics. If logarithmic utility
is introduced, rederive the objective rather than relabeling expected profit.

\paragraph{R5. Establish an observation bridge or restrict the empirical claim. Status: pending.}
Ordinary public quotes are not independent resets. Address initial state,
fundamental-value proxies, endogenous order flow, changing parameters, noise
covariance, sign changes and unobserved tails. A reduced-form VAR score is not
automatically the structural gain in the theorem.
\emph{Completion evidence:} a justified identification/partial-identification
argument under stated conditions, or a forecasting-only interpretation with
certification claims removed from that application. Test the conditions rather
than treating good forecasts as proof of them.

\paragraph{R6. Run a protected real-data evaluation. Status: NOT\_RUN.}
Document event-level quote provenance, instruments, units, dates, fees, carry
and earlier use. Freeze venues/dates and specifications before test access.
Use chronological purging, volatility and OFI/depth baselines, prespecified
inference units and all failure outcomes. An exposure replay must account for
spread, funding, fill/latency assumptions and capital at risk. Distinguish
isolated round trips from a portfolio backtest and from market-wide intervention.
\emph{Completion evidence:} licensed/reproducible data description, immutable
split and code, held-out results after costs, limitations and external replay.
The synthetic unit-test fixture cannot complete R6.

\paragraph{R7. State what happens when certification is unavailable. Status: partial diagnostics.}
The returned bias and future-change cases demonstrate failure. Specify a
defensible response such as abstention or a prior-only restriction; determine
what verified disclosure could restore the missing assumption and what it costs.
Repeated use requires simultaneous or time-uniform control, not repeated use of
a single-decision confidence level.
\emph{Completion evidence:} a stated contract/failure-response rule with a valid
guarantee or an explicit impossibility boundary, tested under prespecified
violations. Larger samples under a wrong mean do not repair the problem.

\paragraph{R8. Obtain independent replication and reconcile claims. Status: pending.}
Have another implementation reproduce the central calculation and meaningful
stress cases. Reconcile abstracts, conclusions, theorem tags, old result tables,
duplicate paper layouts, co-author contributions and AI-use statements. Keep
the original evidence and limitations visible. Treat internal ``proved'' tags
as statements about written arguments until independent review is recorded.
\emph{Completion evidence:} an external replication record, a claim-to-evidence
map, resolved discrepancies and co-author approval for circulation.

\subsection*{Extensions that remain ambitions}
Endogenous dealer depth and quotes, venue-wide gating, equilibrium welfare,
dynamic inventories and AI-adoption externalities are separate extensions.
They should enter the main claim only with a proved mechanism and appropriate
evidence. The new fixed-position experiment does not establish those outcomes.
A strong core contribution need not contain every extension.

\subsection*{Order of work and maintenance rule}
Prioritize R1--R3 together: determine what is new, whether it is true, and when
the information purchase changes an economically meaningful decision. R4--R6
determine the strength of the finance application. R7--R8 determine defensible
use and reproducibility. Professor review can guide and challenge each item;
it is not a substitute for producing the missing proof or evidence.

For each update record the date, task ID, changed claim, immutable code/data
reference, result and remaining limitation. Mark an item completed only when
its named evidence exists. Update this register, its Markdown/JSON source and
all PDF variants together. Publication quality remains a judgment about
originality, correctness and usefulness, not a score produced by the notebook.
\endgroup

% END ADVANCED RETURN STATUS 20261030
\end{document}
